%%%%%%%%%%%%%%%%%%%%%%%%%%%%%%%%%%%%%%%%%%%%%%%%%%%%%%%%%%%%%%%%%%%%%%%%%%%%%%%%%%%%%%%%%
% Tiểu luận/Khóa luận/Luận văn/Luận án tốt nghiệp
% LaTeX Template
% Phiên bản 1.0 (tạo ra ngày 05/10/2022)
% Edited &Modified by VietNC ngày 12/12/2022
%
% Template này có thể được download tại:
% https://github.com/vietnc/HUS-template
% https://github.com/cpc1996/HUS-Dissertation-Template
% 
%
% Phiên bản 1.0 được chỉnh sửa bởi:
% Công Phương Cao (congphuongcao_t59@hus.edu.vn)
% Nguyễn Cảnh Việt (vietncp@gmail.com)

% Template được tham khảo từ một phiên bản template của:
% Steve Gunn (http://users.ecs.soton.ac.uk/srg/softwaretools/document/templates/)
% Sunil Patel (http://www.sunilpatel.co.uk/thesis-template/)
%
%%%%%%%%%%%%%%%%%%%%%%%%%%%%%%%%%%%%%%%%%%%%%%%%%%%%%%%%%%%%%%%%%%%%%%%%%%%%%%%%%%%%%%%%%
%----------------------------------------------------------------------------------------
%	(KHÔNG CHỈNH SỬA PHẦN NÀY)
%
%	PHẦN 1: CÁC PACKAGE CƠ BẢN VÀ CÁC TÙY CHỈNH VĂN BẢN
%----------------------------------------------------------------------------------------
%----------------------------------------------------------------------------------------
%	(KHÔNG CHỈNH SỬA PHẦN NÀY)
%
%	PHẦN 1: CÁC PACKAGE CƠ BẢN VÀ CÁC TÙY CHỈNH VĂN BẢN
%----------------------------------------------------------------------------------------

\documentclass[
12pt,
oneside,
english,
doublespacing,
nolistspacing,
liststotoc,
parskip,
headsepline,
chapterinoneline,
]{MastersDoctoralThesis}

% \usepackage[utf8]{inputenc} 
\usepackage[utf8]{vietnam} 
%\usepackage[T1]{fontenc}

\usepackage{mathptmx}
\usepackage{amsmath}
\allowdisplaybreaks

%https://www.overleaf.com/learn/latex/Biblatex_citation_styles
%\usepackage[backend=bibtex,style=authoryear,natbib=true]{biblatex}
\usepackage[backend=bibtex,style=numeric,citestyle=ieee,natbib=true]{biblatex}

\addbibresource{main.bib}

\usepackage[autostyle=true]{csquotes}


%----------------------------------------------------------------------------------------
%	PHẦN 2: CÁC PACKAGE BỔ TRỢ THÊM VÀO TRONG QUÁ TRÌNH BIÊN SOẠN
%----------------------------------------------------------------------------------------

\RequirePackage{setlst}		% Liệt kê/trích dẫn code

\usepackage{multirow}
\usepackage{subfigure}
\usepackage[fontsize=13pt]{scrextend}

%https://www.sascha-frank.com/latex-font-size.html
%https://tex.stackexchange.com/questions/103286/how-to-change-section-subsection-font-size
\usepackage{titlesec}

\titleformat{\section}
{\normalfont\fontsize{13}{13}\bfseries}{\thesection}{1em}{}

\titleformat{\subsection}
{\normalfont\fontsize{13}{13}\bfseries\itshape}{\thesubsection}{1em}{}

\titleformat{\subsubsection}
{\normalfont\fontsize{13}{13}}{\thesubsubsection}{1em}{}

%https://tex.stackexchange.com/questions/351961/how-to-indent-code-in-beginverbatim
\usepackage{fancyvrb} 		% Fancy Verbatim
\fvset{tabsize=4,vspace=0pt,fontsize=\footnotesize}

\usepackage{longtable} 		% Bảng dài - Long table

\usepackage[figuresright]{rotating} % Bảng ngang - Sideways table
\usepackage{tabularx}

\usepackage{fontawesome5} 	% Các biểu tượng, ký hiệu đặc biệt

\usepackage{tikz} 			% Vẽ hình
\usepackage{indentfirst}
\setlength{\parindent}{0.5cm}	%Thut dau dong cho doan van
\setlength{\parskip}{1.4ex plus 0.5ex minus 0.3ex}		% Tạo khoảng trống giữa 2 đoạn văn %spacing between two paragraph

\usetikzlibrary{calc}
%----------------------------------------------------------------------------------------
%	PHẦN 2: THÔNG TIN VỀ Khóa luận tốt nghiệp/Tiểu luận (THESIS INFORMATION)
% 	Điền thông tin của các bạn trong file(2.thesis-information) này 
%----------------------------------------------------------------------------------------
%----------------------------------------------------------------------------------------
%	PHẦN 3: THÔNG TIN VỀ Khóa luận tốt nghiệp/Tiểu luận (THESIS INFORMATION)
%----------------------------------------------------------------------------------------
\author{NHÓM 6} 			% Ví dụ: "Nguyễn Thị Thu Thảo"
\thesistitle{NHẬN DIỆN CẢM XÚC QUA VĂN BẢN} % Ví dụ: "Nghiên cứu chế tạo và tính chất quang học của vật liệu nano LaF3:Sm3+"

\supervisor {Nguyễn Tiến Cường} 	% Ví dụ: "PGS. TS. Nguyễn Ngọc Long"
\supervisorr{Vi Anh Quân} 	% Tên này nếu không có thì để trống!

\field{Kỹ thuật điện tử và tin học} 					% Ví dụ: "Ngành Vật lý", "Ngành Hóa học", v.v.
\program{Chương trình đào tạo chuẩn} 			% Ví dụ: "Chương trình đào tạo chuẩn", "Chương trình đào tạo cử nhân tài năng", v.v.
\doctype{Báo cáo cuối kỳ} % Điền vào đây: "Khóa luận tốt nghiệp đại học hệ chính quy" hoặc "Tiểu luận"

\university{\href{http://hus.vnu.edu.vn}{TRƯỜNG ĐẠI HỌC KHOA HỌC TỰ NHIÊN}}
\department{\href{http://hus.vnu.edu.vn/gioi-thieu/co-cau-to-chuc/khoa-truc-thuoc/khoa-vat-ly.html}{KHOA VẬT LÝ}}

\AtBeginDocument{
	\hypersetup{pdftitle=\ttitle}
	\hypersetup{pdfauthor=\authorname}
}

%----------------------------------------------------------------------------------------------------------------------------------------
\setcounter{secnumdepth}{4} % Cho phép đánh số định dạng đến cấp thứ 4 (\subsubsection)
\setcounter{tocdepth}{4}    % Nếu muốn cấp này xuất hiện cả trong Mục lục (nhỏ hơn thì tăng lên)
\usepackage{float}
\usepackage{fontawesome5}
\usepackage{twemojis}
\begin{document}
\lstset{style=codeC}	% Thiết lập ngôn ngữ C/C++ là ngôn ngữ mặc định cho phần liệt kê souce code của cả bài
\frontmatter 			% Sử dụng hệ thống đánh số La Mã (i, ii, iii, iv...) cho những trang trước phần mục lục
\pagestyle{plain} 
%----------------------------------------------------------------------------------------
%	(KHÔNG CHỈNH SỬA PHẦN NÀY)%
%	PHẦN 3: TRANG TIÊU ĐỀ/TRANG BÌA (TITLE PAGE)
%----------------------------------------------------------------------------------------
% TRANG BÌA CHÍNH:
\include{3.main-cover}
% TRANG BÌA PHỤ:
\include{4.sub-cover}
%----------------------------------------------------------------------------------------
%	PHẦN 4: DANH NGÔN (QUOTES)
%----------------------------------------------------------------------------------------
%\vspace*{0.2\textheight}

%\noindent\enquote{\itshape 
%	Cái tôi và sự hiểu biết tỷ lệ nghịch với nhau. Hiểu biết càng nhiều cái tôi càng bé. Hiểu biết càng ít, cái tôi càng to.
%}\bigbreak

%\hfill Albert Einstein
%----------------------------------------------------------------------------------------
%	PHẦN 5: LỜI CẢM ƠN (ACKNOWLEDGEMENTS)
%	Các bạn ghi lời cảm ơn vào file(5.ack) này.
%----------------------------------------------------------------------------------------
%----------------------------------------------------------------------------------------
%	PHẦN 6: LỜI CẢM ƠN (ACKNOWLEDGEMENTS)
%----------------------------------------------------------------------------------------

\begin{acknowledgements}
	\addchaptertocentry{\acknowledgementname}
	\thispagestyle{empty}
	
	Lời đầu tiên, nhóm chúng em xin gửi lời cảm ơn chân thành và sâu sắc nhất đến các thầy cô giáo Trường Đại học Khoa học Tự nhiên, Đại học Quốc gia Hà Nội. Trong suốt quá trình học tập dưới mái trường, các thầy cô đã tận tình truyền dạy và trang bị cho nhóm những kiến thức nền tảng vô cùng quý báu về trí tuệ nhân tạo, đại số tuyến tính, giải tích,..., tạo tiền đề vững chắc để em thực hiện đề tài này.
	
	Đặc biệt, nhóm xin bày tỏ lòng biết ơn sâu sắc đến TS. Nguyễn Tiến Cường và CN. Vi Anh Quân. Sự hỗ trợ chuyên môn và những kiến thức bổ trợ quý báu về lý thuyết Học máy của hai thầy chính là kim chỉ nam giúp nhóm chúng em giải quyết bài toán “Xây dựng hệ thống nhận diện đa nhãn cảm xúc”, từ đó hoàn thiện nghiên cứu một cách bài bản, khoa học và hiệu quả hơn.
	
	Nhóm cũng xin cảm ơn các bạn cùng lớp Khóa 68 Kỹ thuật Điện tử và Tin học đã luôn chia sẻ tài liệu, đóng góp ý kiến phản biện và cổ vũ tinh thần cho nhóm trong suốt thời gian triển khai đề tài cuối kỳ môn Học Máy.
	
	Mặc dù đã có nhiều cố gắng để hoàn thiện đề tài một cách chỉn chu nhất, song do kiến thức và kinh nghiệm thực tế của bản thân còn hạn chế, quyển báo cáo chắc chắn không tránh khỏi những thiếu sót trong quá trình trình bày và thiết kế. Nhóm rất mong nhận được những ý kiến đóng góp và nhận xét từ quý thầy cô để có thể tiếp tục hoàn thiện đề tài cũng như nâng cao kiến thức và kỹ năng của bản thân trong thời gian tới.
	
	\begin{flushright}
		\begin{tabular}{r}
			Xin chân thành cảm ơn! \\[0.6cm]
			\textit{TP. Hà Nội, ngày 03 tháng 06 năm 2026} \\[0.2cm]
			\textbf{Nhóm sinh viên thực hiện} \\[0.6cm]
			Nhóm 6
		\end{tabular}
	\end{flushright}
\end{acknowledgements}
%----------------------------------------------------------------------------------------
%	(KHÔNG CHỈNH SỬA PHẦN NÀY)
%
%	PHẦN 6: MỤC LỤC (LIST OF CONTENTS/FIGURES/TABLES PAGES)
%	ĐẶT KHOẢNG CÁCH GIŨA CÁC TIÊU ĐỀ CỦA PHẦN 7
%----------------------------------------------------------------------------------------
\begin{spacing}{1.15}
	\tableofcontents 	% In ra mục lục chính
\end{spacing}

\begin{spacing}{1.15}
	\listoffigures 		% In ra danh sách hình vẽ
\end{spacing}

\begin{spacing}{1.15}
	\listoftables		% In ra danh sách bảng
\end{spacing}

%----------------------------------------------------------------------------------------
%	PHẦN 7: DANH SÁCH TÊN VIẾT TẮT (ABBREVIATIONS)
%
%	Các từ viết tắt các bạn ghi ở đây
%----------------------------------------------------------------------------------------
\begin{abbreviations}{ll} % Thêm danh sách tên viết tắt (dưới dạng một bảng có 2 cột)

\textbf{NLP} & \textbf{N}atural \textbf{L}anguage \textbf{P}rocessing\\
\textbf{BoW} & \textbf{B}ag of \textbf{W}ord\\
\textbf{TF - IDF} & \textbf{T}erm \textbf{F}requency -  \textbf{I}nvert \textbf{D}ocument \textbf{F}requency\\
\textbf{SVM} & \textbf{S}upport \textbf{V}ector \textbf{M}achine\\
\textbf{RF} & \textbf{R}andom \textbf{F}orest\\
\textbf{CNN} & \textbf{C}onvolution \textbf{N}eural \textbf{N}etwork\\
\textbf{LSTM} & \textbf{L}ong \textbf{S}hort - \textbf{T}erm \textbf{M}emory\\

\end{abbreviations}
%----------------------------------------------------------------------------------------
%	PHẦN 8: CÁC HẰNG SỐ VẬT LÝ/THÔNG SỐ KĨ THUẬT (PHYSICAL CONSTANTS/OTHER DEFINITIONS)
%----------------------------------------------------------------------------------------
%\begin{constants}{lr@{${}={}$}l} % Thêm danh sách các hằng số (dưới dạng một bảng có 3 cột)
%
%%	Lệnh \SI{}{} được cung cấp bởi gói siunitx, hãy đọc tài liệu hướng dẫn để biết cách sử dụng nó
%
%%	Tên hằng số 	 & $Biểu tượng$	& $Hằng số$ cùng với đơn vị\\
%	Vận tốc ánh sáng & $c_{0}$		& \SI{2.99792458e8}{\meter\per\second} (chính xác)\\
%
%\end{constants}
%----------------------------------------------------------------------------------------
%	PHẦN 9: DANH SÁCH KÝ HIỆU (SYMBOLS)
%----------------------------------------------------------------------------------------

%\begin{symbols}{lll} % Thêm danh sách các ký hiệu (dưới dạng một bảng có 3 cột)
%
%%   Ký hiệu	& Ý nghĩa		& Đơn vị \\
%	$a$		& khoảng cách	& \si{\meter} \\
%	$P$		& công suất		& \si{\watt} (\si{\joule\per\second}) \\
%
%	\addlinespace % Khoảng cách để phân biệt giữa ký hiệu Latin với ký hiệu La Mã
%
%	$\omega$ & tần số góc	& \si{\radian} \\
%
%\end{symbols}

%----------------------------------------------------------------------------------------
%	(KHÔNG CHỈNH SỬA PHẦN NÀY)
%
%	PHẦN 10: LỜI ĐỀ TẶNG (DEDICATION)
%----------------------------------------------------------------------------------------

%\dedicatory{Dành tặng/Dành cho/Gửi tới\ldots} 


%----------------------------------------------------------------------------------------
%	PHẦN 11: NỘI DUNG/CÁC CHƯƠNG Khóa luận tốt nghiệp/Tiểu luận (THESIS CONTENT - CHAPTERS)
%----------------------------------------------------------------------------------------

\mainmatter % Bắt đầu đánh số trang (1,2,3...)

\pagestyle{plain}

% Hãy thêm những chương (chapter) của khóa luận/tiểu luận vào thư mục Chapters
% Hãy bỏ chú thích những dòng nếu bạn đã bổ sung những chương vào

% MỞ ĐẦU (LỜI MỞ ĐẦU - Chương 0)

\chapter*{MỞ ĐẦU} % Tên của chương
\addcontentsline{toc}{chapter}{MỞ ĐẦU} % Thêm tên chương vào mục lục

\label{Chapter0} % Để trích dẫn chương này ở chỗ nào đó trong bài, hãy sử dụng lệnh \ref{Chapter0} 

%----------------------------------------------------------------------------------------

Trong thời đại số, lượng dữ liệu văn bản khổng lồ do người dùng sinh ra trên các nền tảng mạng xã hội và diễn đàn ẩn chứa những thông tin giá trị về trạng thái tâm lý của chính họ. Tuy nhiên, các bài toán phân loại cảm xúc truyền thống (chỉ đánh giá đơn nhãn hoặc chia ba mức độ: tích cực, tiêu cực, trung tính) chưa đủ sức phản ánh sự phức tạp trong nội tâm con người, khi một câu nói có thể chứa nhiều sắc thái cùng lúc. Vì vậy, đề tài này được thực hiện với hai mục tiêu cốt lõi: (1) nghiên cứu, xây dựng hệ thống phân loại đa nhãn có khả năng nhận diện 28 trạng thái cảm xúc khác nhau của người dùng từ văn bản và (2) triển khai thực nghiệm, đánh giá, so sánh hiệu quả của các thuật toán học máy trong xử lý ngôn ngữ tự nhiên.
 
Để đảm bảo tính khoa học và tập trung, đề tài được giới hạn trong các đối tượng và phạm vi cụ thể. Về mặt dữ liệu, nghiên cứu sử dụng bộ dữ liệu chuẩn GoEmotions do Google cung cấp, tập trung phân tích các bình luận tiếng Anh. Về phương pháp biểu diễn ngôn ngữ, hệ thống tập trung khai thác và so sánh hai kỹ thuật trích xuất đặc trưng văn bản bao gồm phương pháp thống kê truyền thống TF-IDF và kỹ thuật nhúng từ (Word Embedding). Về phương pháp huấn luyện, đề tài tiến hành đánh giá hiệu năng trên tập hợp các thuật toán bao gồm: Logistic Regression, Random Forest, Linear SVC,  Mạng nơ-ron tích chập (CNN) và Mạng bộ nhớ dài - ngắn (LSTM).

Để làm rõ hơn về đề tài này, bài báo cáo được chia làm các nội dung như sau:

\textbf{Chương 1: Tổng quan.}

\textbf{Chương 2: Các cơ sở lý thuyết.}

\textbf{Chương 3: Thực nghiệm.}

\textbf{Chuong 4: Kết quả.}

\textbf{Chương 5: Kết luận.}
% Chương 1

\chapter{TỔNG QUAN} % Tên của chương

\label{Chapter1} % Để trích dẫn chương này ở chỗ nào đó trong bài, hãy sử dụng lệnh \ref{Chapter1} 

%----------------------------------------------------------------------------------------

% Định nghĩa một số lệnh cần thiết để điều chỉnh định dạng cho một số nội dung nhất định trong bài
\newcommand{\keyword}[1]{\textbf{#1}}
\newcommand{\tabhead}[1]{\textbf{#1}}
\newcommand{\code}[1]{\texttt{#1}}
\newcommand{\file}[1]{\texttt{\bfseries#1}}
\newcommand{\option}[1]{\texttt{\itshape#1}}

%----------------------------------------------------------------------------------------

\section{Ngôn ngữ tự nhiên}

Ngôn ngữ tự nhiên là công cụ biểu đạt và giao tiếp cốt lõi được con người hình thành, sử dụng để tương tác hàng ngày trong xã hội, chẳng hạn như tiếng Việt, tiếng Anh, tiếng Trung... Theo các thống kê định lượng trong ngôn ngữ học máy tính, trên thế giới hiện nay có khoảng hơn 5.600 ngôn ngữ khác nhau được phân bố không đồng đều và chỉ có một tỷ lệ nhỏ trong số đó sở hữu hệ thống chữ viết hoàn chỉnh \cite{jurafsky2020speech}. Khác biệt hoàn toàn với các hệ thống ngôn ngữ lập trình hay ngôn ngữ hình thức vốn vận hành dựa trên các quy tắc toán học bất biến, ngôn ngữ tự nhiên sở hữu cấu trúc hình thái phức tạp, linh hoạt và luôn chứa đựng nhiều sự mơ hồ \cite{singh2018natural}.

Các đặc điểm hình thái và cú pháp cốt lõi của ngôn ngữ tự nhiên bao gồm:
\begin{itemize}
\item \textbf{Tính biến đổi}: Ngôn ngữ không mang trạng thái cố định mà liên tục biến đổi, tiến hóa theo thời gian, văn hóa và bối cảnh lịch sử của cộng đồng sử dụng \cite{singh2018natural}.

\item \textbf{Tính đa nghĩa}: Một từ hoặc một cấu trúc câu trong ngôn ngữ tự nhiên có thể mang nhiều tầng ý nghĩa khác nhau tùy thuộc vào ngữ cảnh xuất hiện \cite{singh2018natural}. Ví dụ, từ "bank" trong tiếng Anh có thể đại diện cho tổ chức tín dụng ngân hàng hoặc bờ sông, đòi hỏi máy tính phải có khả năng khử nhập nhằng ngữ nghĩa \cite{singh2018natural}.

\item \textbf{Tính ngữ cảnh}: Ngữ nghĩa thực tế của một thông điệp phụ thuộc chặt chẽ vào tình huống giao tiếp, đặc điểm nhân khẩu học của người nói, người nghe cũng như môi trường xung quanh \cite{singh2018natural}.

\item \textbf{Tính không hoàn hảo}: Văn bản ngôn ngữ tự nhiên, đặc biệt là dữ liệu dạng văn bản thô trích xuất từ các nền tảng mạng xã hội, thường có mật độ nhiễu rất cao do chứa nhiều lỗi ngữ pháp, lỗi viết tắt, từ lóng hoặc lỗi chính tả cấu trúc câu không hoàn chỉnh \cite{singh2018natural}.

\item \textbf{Tính hệ thống}: Mặc dù biến thiên liên tục, ngôn ngữ tự nhiên vẫn tuân thủ các quy tắc tổ chức cú pháp nghiêm ngặt ở cấp độ câu để đảm bảo thông điệp truyền tải có thể dự đoán và hiểu được \cite{singh2018natural}.

\item \textbf{Tính đa phương thức}: Quá trình truyền tải thông điệp không chỉ dựa vào ký tự văn bản thuần túy mà còn kết hợp chặt chẽ với ngữ điệu, âm lượng hoặc các yếu tố phi ngôn từ để biểu đạt sắc thái biểu cảm \cite{singh2018natural}.

\item \textbf{Tính phân cấp}: Cấu trúc của ngôn ngữ tự nhiên được tổ chức phân cấp một cách chặt chẽ từ các đơn vị nhỏ nhất là âm vị, hình vị, từ vựng cho đến các cấu trúc vĩ mô hơn như cụm từ, câu, đoạn văn và toàn bộ văn bản \cite{singh2018natural}.
\end{itemize}

Chính những đặc tính phức tạp và mơ hồ này đã tạo ra rào cản kỹ thuật rất lớn cho máy tính, biến việc xử lý dữ liệu chữ thô trở thành một tác vụ tốn kém và đòi hỏi nhiều lao động thủ công \cite{singh2018natural}. Việc thấu hiểu các đặc điểm này là tiền đề bắt buộc để xây dựng các mô hình học máy phân loại cảm xúc tự động ở các chương sau \cite{hastie2009elements},\cite{hochreiter1997long}.

%----------------------------------------------------------------------------------------

\section{Ngôn ngữ Tiếng Anh}

Tiếng Anh thuộc loại hình ngôn ngữ phân tích, đồng thời vẫn lưu giữ một số đặc tính của ngôn ngữ hòa kết. Điều này đồng nghĩa với việc các mối quan hệ cú pháp và ngữ pháp trong câu được biểu thị chủ yếu thông qua trật tự từ và các hư từ (bao gồm giới từ, trợ động từ) thay vì biến đổi hình thái từ quá phức tạp như các ngôn ngữ thuần hòa kết. Tuy nhiên, khác với các ngôn ngữ đơn lập thuần túy như tiếng Việt, từ vựng tiếng Anh có cấu tạo dựa trên sự kết hợp chặt chẽ của các hình vị bao gồm gốc từ, tiền tố và hậu tố thành một khối thống nhất.

Dựa trên các nghiên cứu ngôn ngữ học định lượng về phân tích máy tính của Kučera và Francis (1967), sự phân bố của hơn 170.000 từ vựng tiếng Anh thông dụng theo cấu trúc số lượng âm tiết được xác định cụ thể như sau \cite{kucera1967computational}:
\begin{itemize}
\item Khoảng 15\% số từ trong từ điển là từ đơn âm tiết (tuy nhiên nhóm từ này lại chiếm tần suất xuất hiện áp đảo, chiếm tới hơn 40\% trong văn bản và giao tiếp thực tế hàng ngày) \cite{kucera1967computational}.

\item Khoảng 35\% số từ là từ sở hữu hai âm tiết \cite{kucera1967computational}.

\item Khoảng 30\% số từ là từ sở hữu ba âm tiết \cite{kucera1967computational}.

\item Khoảng 20\% còn lại là các từ đa âm tiết (có cấu trúc từ bốn âm tiết trở lên) \cite{kucera1967computational}.
\end{itemize}

Chính vì đặc thù của loại hình ngôn ngữ phân tích, việc sắp xếp các từ trong tiếng Anh theo một trật tự nghiêm ngặt là yếu tố cốt lõi để phân định vai trò ngữ pháp và ngữ nghĩa của câu. Ví dụ, cấu trúc câu "\textit{The dog bit the man}" (Con chó cắn người đàn ông) mang ý nghĩa đối lập hoàn toàn với "\textit{The man bit the dog}" (Người đàn ông cắn con chó). Nhờ vào kết cấu vị trí này, hệ thống mới phân định rõ ràng đâu là chủ thể hành động và đâu là đối tượng chịu tác động. Trong tiếng Anh, trật tự kết cấu câu chuẩn mực và phổ biến nhất tuân theo sơ đồ \textbf{SVO} (Chủ ngữ - Động từ - Tân ngữ: Subject - Verb - Object). Đặc điểm hệ thống về mặt hình thái và cú pháp này là cơ sở quan trọng để các công cụ tiền xử lý thực hiện bóc tách, đưa từ về dạng nguyên thể, từ đó giúp các mô hình học máy trích xuất chính xác mối quan hệ ngữ cảnh và phân loại đa nhãn cảm xúc \cite{singh2018natural}.

\section{Xử lý ngôn ngữ tự nhiên}

Xử lý ngôn ngữ tự nhiên (Natural Language Processing - NLP) là một lĩnh vực giao thoa cốt lõi giữa Trí tuệ nhân tạo (AI) và Khoa học máy tính, tập trung vào việc nghiên cứu các kỹ thuật cho phép máy tính có khả năng đọc, hiểu, diễn giải và tạo ra ngôn ngữ của con người một cách tự động và có ý nghĩa \cite{singh2018natural}. Về mặt phân cấp hệ thống, NLP không hoạt động độc lập mà có mối quan hệ giao thoa chặt chẽ với Học máy (ML) và Học sâu (DL) trong bức tranh tổng thể của Trí tuệ nhân tạo (AI), được mô tả trực quan qua sơ đồ ở Hình 1.1. 
%\begin{figure}[h!]
%	\centering
%	\includegraphics[width=0.8\textwidth]{VNU Dai hoc Tong hop Ha Noi - 19 Le Thanh Tong.jpeg}
%	\caption[Đại học Tổng hợp]{Đại học Tổng hợp Hà Nội tại 19 Lê Thánh Tông.}
%	\label{fig:DHTHHN}
%\end{figure}
\begin{figure}[htbp]
	\centering
	\includegraphics[width=0.8\textwidth]{1.1.png}
	\caption{Sơ đồ vị trí và mối quan hệ giao thoa của NLP trong hệ sinh thái AI, ML và DL.}
	\label{fig:bieudo1}
\end{figure}

Sự giao thoa được thể hiện rõ nét qua việc các mô hình NLP hiện đại ngày càng dịch chuyển từ hệ thống dựa trên luật cứng nhắc sang việc áp dụng các thuật toán Học máy để tự động khai phá thông tin \cite{singh2018natural}. Ở tầng sâu hơn, các kiến trúc Học sâu đóng vai trò như các bộ lọc mạnh mẽ giúp trích xuất và tối ưu hóa các đặc trưng ngôn ngữ phức tạp mà các phương pháp truyền thống không thể chạm tới \cite{hastie2009elements},\cite{hochreiter1997long} 
Để hiện thực hóa các mô hình công nghệ này vào thực tế bài toán, đề tài tuân thủ nghiêm ngặt quy trình phát triển một hệ thống NLP chuẩn mực công nghiệp. Quy trình này vận hành theo một vòng lặp khép kín gồm 6 giai đoạn cốt lõi, chia thành 3 phân vùng trách nhiệm lớn như minh họa tại Hình 1.2: 
\begin{figure}[htbp]
	\centering
	\includegraphics[width=0.8\textwidth]{1.2.png}
	\caption{Sơ đồ chu kỳ phát triển toàn diện hệ thống NLP của dự án.}
	\label{fig:bieudo2}
\end{figure}

Dựa trên cấu trúc từ Hình 1.2, quy trình phát triển tổng quát điều phối hệ thống qua các phân vùng chức năng sau:
\begin{itemize}
\item \textbf{Phân vùng Nghiên cứu (Research)}: Tiến hành thu thập tập dữ liệu văn bản thô (Data collection) , sau đó thực hiện phân tích đặc trưng ngôn ngữ, làm sạch, chuẩn hóa cấu trúc dữ liệu thô và thực nghiệm kiểm thử thuật toán (Analysis \& experimenting).

\item \textbf{Phân vùng Phát triển (Development)}: Đưa ma trận đặc trưng văn bản vào huấn luyện các thuật toán phân lớp cốt lõi (Training) và tiến hành viết mã nguồn, tối ưu hóa các hàm toán học để cấu hình hệ thống (Implementation).

\item \textbf{Phân vùng Vận hành (Operations)}: Đóng gói và triển khai mô hình hoàn chỉnh lên môi trường ứng dụng thực tế (Deploying) , đồng thời thiết lập hệ thống liên tục giám sát sai số, theo dõi độ ổn định của dữ liệu người dùng theo thời gian (Monitoring).
\end{itemize}

Sự phối hợp nhịp nhàng của chu kỳ phát triển tổng quát này là nền tảng kỹ thuật quyết định để đề tài tiến tới giải quyết bài toán nhận diện sắc thái tinh vi trong văn bản mạng xã hội ở các mục tiếp theo.

\section{Nhận diện cảm xúc trong văn bản}

\subsection{Khái niệm về cảm xúc dưới góc nhìn khoa học}

Trước khi tiếp cận dưới góc độ công nghệ máy tính, cần làm rõ bản chất hệ thống: Cảm xúc là một trạng thái phản ứng sinh lý và tâm lý phức tạp của con người trước các tác nhân kích thích từ môi trường bên ngoài hoặc từ các luồng suy nghĩ nội tại \cite{jurafsky2020speech}. Theo học thuyết tâm lý học kinh điển của Paul Ekman, con người sở hữu một tập hợp gồm 6 trạng thái cảm xúc nền tảng mang tính phổ quát toàn cầu bao gồm: \textit{Vui vẻ (Joy), Buồn bã (Sadness), Giận dữ (Anger), Sợ hãi (Fear), Ngạc nhiên (Surprise), và E ngại/E dè (Disgust)}, được trực quan hóa thông qua hệ thống biểu cảm tương đương tại Hình 1.3 \cite{jurafsky2020speech}
\begin{figure}[htbp]
	\centering
	
\includegraphics[width=0.8\textwidth]{1.3.png}
	\caption{Minh họa hệ thống 6 trạng thái cảm xúc nền tảng theo học thuyết Paul Ekman.}
	\label{fig:bieudo3}
\end{figure}

Trong tự nhiên, trạng thái tâm lý này của con người được phát tín hiệu và biểu đạt ra thế giới bên ngoài thông qua ba phương thức cốt lõi \cite{jurafsky2020speech}:
\begin{enumerate}
	\item \textbf{Phương thức trực quan}: Biểu cảm qua hệ thống cơ mặt (nụ cười, ánh mắt, cái nhíu mày) và các cử chỉ, hành vi của cơ thể.
	
	\item \textbf{Phương thức thính giác}: Biểu cảm qua giọng nói, thể hiện bằng các biến thiên về mặt tần số, cao độ, âm lượng, tốc độ và ngữ điệu phát âm.
	
	\item \textbf{Phương thức ngôn từ}: Biểu cảm qua văn bản chữ viết, được thể hiện một cách tinh vi thông qua việc lựa chọn từ vựng sắc thái, cấu trúc cú pháp câu, các kỹ thuật châm biếm, hoán dụ, hoặc hệ thống dấu biểu cảm (!, ?) và từ lóng trên các nền tảng số.
\end{enumerate}

Trong thực tế tương tác hàng ngày và đặc biệt là trên không gian mạng xã hội, các trạng thái tâm lý này không tồn tại đơn lẻ mà liên tục biến đổi, đan xen và mở rộng thành các sắc thái biểu cảm phức tạp hơn (ví dụ: từ \textit{Vui vẻ} phát triển thành \textit{Ngưỡng mộ - Admiration}, \textit{Thích thú - Amusement}, hay \textit{Lòng biết ơn - Gratitude}) \cite{jurafsky2020speech}. Do giới hạn về mặt thu thập dữ liệu kỹ thuật số, đề tài này tập trung khai thác sâu sắc vào phương thức ngôn từ để làm tiền đề cho máy tính tự động trích xuất và phân tích tri thức cảm xúc từ văn bản của người dùng \cite{jurafsky2020speech}. 

\subsection{Bài toán nhận diện cảm xúc từ văn bản}

Nhận diện cảm xúc trong văn bản là một nhánh phát triển chuyên sâu và phức tạp của bài toán Phân tích quan điểm và Khai phá ý kiến trong lĩnh vực xử lý ngôn ngữ tự nhiên \cite{singh2018natural}. Trong khi phân tích quan điểm truyền thống chỉ dừng lại ở việc phân loại tính phân cực tài liệu theo các hướng thô như Tích cực (Positive), Tiêu cực (Negative), hoặc Trung tính (Neutral), thì nhận diện cảm xúc hướng tới việc bóc tách các trạng thái tâm lý tinh vi, có chiều sâu hơn của con người (ví dụ: giận dữ, biết ơn, lo lắng, hay ngưỡng mộ).

Dữ liệu văn bản trên mạng xã hội và các diễn đàn trực tuyến như Reddit phản ánh chân thực trạng thái tâm lý thời gian thực của người dùng, tạo ra nguồn tài nguyên khép kín cho các bài toán phân tích xã hội và chăm sóc khách hàng \cite{singh2018natural}. Tuy nhiên, việc máy tính tự động nhận diện cảm xúc từ các chuỗi ký tự thô đối mặt với những thách thức học máy rất lớn \cite{singh2018natural}:
\begin{itemize}
\item \textbf{Sự mơ hồ và ẩn dụ của ngôn từ}: Con người thường không biểu lộ cảm xúc một cách trực diện qua các từ khóa cố định mà lồng ghép qua các cấu trúc châm biếm, hoán dụ, hoặc ẩn dụ ngữ cảnh \cite{singh2018natural}.

\item \textbf{Đặc thù ngôn ngữ mạng xã hội}: Các văn bản thu thập được thường rất ngắn, cấu trúc ngữ pháp lỏng lẻo, viết tắt tự do và thay đổi liên tục theo xu hướng, làm suy giảm nghiêm trọng hiệu năng của các bộ trích xuất đặc trưng truyền thống \cite{singh2018natural}.
\end{itemize}

Chính vì vậy, việc chuyển hóa các chuỗi chữ thô mang sắc thái biểu cảm thành các không gian vectơ toán học, sau đó áp dụng các thuật toán học máy và học sâu để phân lớp ngữ cảnh là chìa khóa cốt lõi để máy tính có thể "thấu cảm" được văn bản của con người \cite{hastie2009elements},\cite{singh2018natural},\cite{hochreiter1997long}.

\section{Bài toán phân loại đa nhãn}

Trong các bài toán học máy phân loại truyền thống (Single-label Classification), mỗi mẫu dữ liệu đầu vào chỉ được phép ánh xạ duy nhất tới một nhãn lớp độc nhất loại trừ lẫn nhau (ví dụ: phân loại đa lớp tích cực, tiêu cực hoặc trung tính). Tuy nhiên, ngôn ngữ trên mạng xã hội phản ánh nội tâm phức tạp của con người, nơi một câu nói có thể chứa đựng đồng thời nhiều trạng thái tâm lý đan xen \cite{demszky2020goemotions}.

\begin{figure}[H]
	\centering
	\includegraphics[width=0.8\textwidth]{1.5.png}
	\caption{Ví dụ minh họa về văn bản một nhãn và đa nhãn cảm xúc.}
	\label{fig:bieudo1.5}
\end{figure}

Do đó, đề tài này tiếp cận Bài toán Phân loại Đa nhãn (Multi-label Classification) \cite{demszky2020goemotions}. Về mặt toán học, cho trước một không gian văn bản đầu vào $x$ và một tập hợp gồm $L$ nhãn cảm xúc độc lập $Y = \{Y_1, Y_2, \dots, Y_L\}$. Mục tiêu của mô hình là học một hàm quyết định $h(x)$ để dự đoán ra một vectơ đầu ra dạng nhị phân $\hat{y} = [\hat{y_1}, \hat{y_2}, \dots, \hat{y_L}] \in \{0, 1\}^L$ trong đó $\hat{y_j} = 1$ nếu mô hình nhận diện văn bản chứa sắc thái cảm xúc thứ $j$, và bằng 0 trong trường hợp ngược lại. Lúc này, các nhãn cảm xúc không còn loại trừ lẫn nhau, cho phép một bình luận có thể được gán đồng thời nhiều nhãn như vừa "\textit{admiration}" (ngưỡng mộ) vừa "\textit{joy}" (vui sướng) \cite{demszky2020goemotions}.

%----------------------------------------------------------------------------------------

\section{Kiến trúc tổng quan của dự án}

Để giải quyết bài toán phân loại đa nhãn cảm xúc từ văn bản mạng xã hội, đề tài xây dựng và triển khai hệ thống dựa trên một quy trình xử lý dữ liệu học máy chuẩn mực. Kiến trúc tổng quan của dự án bao gồm các khối chức năng cốt lõi được điều phối tuần tự theo sơ đồ minh họa ở Hình 1.5.
\begin{figure}[H]
	\centering
	\includegraphics[width=0.8\textwidth]{1.4.png}
	\caption{Sơ đồ kiến trúc luồng xử lý dữ liệu và huấn luyện hệ thống.}
	\label{fig:bieudo4}
\end{figure}

Dựa trên cấu trúc mô hình hóa từ Hình 1.5, luồng vận hành chi tiết của dự án được đặc tả qua 5 giai đoạn chính sau:
\begin{itemize}
\item Tập dữ liệu và Tiền xử lý dữ liệu: Hệ thống tiếp nhận tập dữ liệu thô GoEmotions gồm các bình luận từ nền tảng Reddit \cite{demszky2020goemotions}. Do đặc thù ngôn ngữ mạng xã hội chứa nhiều nhiễu, các chuỗi văn bản sẽ được đi qua hàm tiền xử lý nâng cao nhằm chuẩn hóa bảng mã, khai triển từ viết tắt, xử lý cấu trúc phủ định và lemmatize đưa từ về dạng nguyên thể.

\item Trích xuất đặc trưng: Các văn bản sau khi làm sạch hoàn toàn sẽ được số hóa thành các vectơ toán học. Giai đoạn này thực hiện thực nghiệm và trích xuất đặc trưng dựa trên phương pháp thống kê TF-IDF kết hợp với các kỹ thuật nhúng từ không gian liên tục như Word2Vec và FastText \cite{singh2018natural},\cite{mikolov2013efficient}.

\item Phân chia dữ liệu thực nghiệm: Sau khi ma trận đặc trưng được hình thành, dữ liệu sẽ được phân tách độc lập thành hai tập con:
\begin{itemize}
\item[$\circ$] Dữ liệu huấn luyện (Training set): Chiếm tỷ lệ lớn (ví dụ: 80\%), phục vụ cho quá trình học các đặc trưng cảm xúc của các thuật toán.

\item[$\circ$] Dữ liệu đánh giá (Validation/Test set): Chiếm tỷ lệ còn lại (ví dụ: 20\%), được giữ hoàn toàn độc lập để kiểm thử và đánh giá khách quan năng lực phân loại của hệ thống.
\end{itemize}
\item Huấn luyện mô hình: Tập dữ liệu huấn luyện được đưa vào các thuật toán phân loại cốt lõi (bao gồm Logistic Regression, Random Forest, LinearSVC, CNN, LSTM) để tối ưu hóa hệ thống trọng số. Quá trình này giúp chuyển đổi các thuật toán từ dạng "Mô hình" lý thuyết ban đầu trở thành một "Mô hình đã huấn luyện" mang đầy đủ tri thức phân lớp.

\item Phân loại đầu ra: Mô hình đã huấn luyện sẽ tiếp nhận tập dữ liệu đánh giá, tính toán và đưa ra kết quả dự đoán nhãn cảm xúc dưới dạng một vectơ nhị phân đa nhãn, hoàn thành mục tiêu nhận diện đồng thời các trạng thái tâm lý đan xen trong văn bản người dùng.
\end{itemize}



% Chương 2

\chapter{CƠ SỞ LÝ THUYẾT} % Tên của chương

\label{Chapter2} % Để trích dẫn chương này ở chỗ nào đó trong bài, hãy sử dụng lệnh \ref{Chapter2} 

%----------------------------------------------------------------------------------------

\section{Các phương pháp biểu diễn văn bản}

Trong lĩnh vực xử lý ngôn ngữ tự nhiên (NLP), máy tính không thể trực tiếp tiếp nhận và tính toán trên các chuỗi ký tự hay từ vựng thô dạng văn bản \cite{singh2018natural}. Do đó, biểu diễn văn bản là quá trình ánh xạ các đơn vị ngôn ngữ (từ, cụm từ hoặc toàn bộ văn bản) thành các vectơ số thực trong không gian toán học \cite{singh2018natural}. Hiệu năng của các mô hình phân loại cảm xúc phụ thuộc rất lớn vào khả năng lưu giữ thông tin ngữ nghĩa và cú pháp của phương pháp biểu diễn dữ liệu đầu vào này \cite{singh2018natural},\cite{mikolov2013efficient}.

\subsection{Phương pháp TF - IDF}

TF-IDF là một kĩ thuật phổ biến được sử dụng trong nhiều bài toán NLP và khai phá dữ liệu dạng văn bản nhằm mục đích tính toán trọng số (độ quan trọng) của một từ trong một văn bản cụ thể, khi văn bản đó nằm trong một tập hợp nhiều văn bản khác nhau (corpus) \cite{singh2018natural}.

TF-IDF giúp chuyển đổi dạng biểu diễn văn bản thô thành dạng không gian vectơ (VSM) hoặc thành những vectơ thưa thớt (sparse vectors) phục vụ cho các mô hình tính toán \cite{singh2018natural}. Kỹ thuật này giúp làm tăng giá trị trọng số tỷ lệ thuận với số lần xuất hiện của từ trong văn bản, nhưng lại tỷ lệ nghịch với tần suất từ đó xuất hiện trên toàn bộ tập dữ liệu đầu vào \cite{singh2018natural}. 

Kỹ thuật này được cấu thành từ hai chỉ số thành phần:
\begin{itemize}
\item \textbf{TF (Term Frequency)}: Tần suất xuất hiện của một từ trong một đoạn văn bản. Khi tính toán độc lập, mức độ quan trọng của các từ là như nhau \cite{singh2018natural}. 

\item \textbf{IDF (Inverse Document Frequency)}: Được dùng để đánh giá mức độ quan trọng của một từ trên toàn tập dữ liệu. Do trong văn bản thường xuất hiện nhiều từ không quan trọng xuất hiện với tần suất cao (ví dụ: stop words), IDF được sử dụng để làm giảm đi mức độ quan trọng của các từ đó \cite{singh2018natural}. 
\end{itemize}
Công thức xác định các thành phần của TF-IDF được cho bởi:
\begin{equation}
	TF_i = \frac{n_i}{N_i}
\end{equation}

Trong đó:
\begin{itemize}
	\item $i$: $1 \dots D$.
	\item $n_i$: Tần số xuất hiện của từ trong văn bản $i$.
	\item $N_i$: Tổng số từ trong văn bản $i$.
\end{itemize}
\begin{equation}
	IDF_i = \log(\frac{D}{d})
\end{equation}

Trong đó:
\begin{itemize}
	\item D: Tổng số document trong tập dữ liệu
	\item d: Số lượng document có sự xuất hiện của từ
\end{itemize}

Khi đó TD - IDF tổng hợp cuối cùng:
\begin{equation}
	TF_i - IDF_i = TF_i\times{IDF_i}
\end{equation}

\subsection{Phương pháp Word2Vec}

Word2Vec là một kỹ thuật biểu diễn từ vựng dưới dạng các vectơ số thực trong không gian liên tục, được đề xuất bởi nhóm nghiên cứu Google Brain do Mikolov và cộng sự phát triển năm 2013 \cite{mikolov2013efficient}. Khác với các phương pháp biểu diễn rời rạc như TF-IDF, Word2Vec ánh xạ các từ có ngữ nghĩa tương đồng tới các vị trí gần nhau trong không gian vectơ nhiều chiều, từ đó cho phép mô hình học máy khai thác được mối quan hệ ngữ nghĩa giữa các từ \cite{mikolov2013efficient}.

Ý tưởng cốt lõi của Word2Vec dựa trên giả thuyết phân bố trong ngôn ngữ học, cho rằng những từ xuất hiện trong các ngữ cảnh tương tự thường mang ý nghĩa tương tự nhau \cite{mikolov2013efficient}. Thông qua quá trình huấn luyện trên một tập văn bản lớn, mô hình học được các biểu diễn vectơ có khả năng phản ánh cả đặc trưng ngữ nghĩa và cú pháp của từ vựng \cite{mikolov2013efficient}.

Word2Vec sử dụng mạng nơ-ron nông gồm một lớp ẩn để học biểu diễn từ và được triển khai theo hai kiến trúc chính \cite{mikolov2013efficient}:
\begin{itemize}
\item \textbf{Continuous Bag-of-Words (CBOW)}: Mô hình sử dụng các từ ngữ cảnh xung quanh để dự đoán từ trung tâm. Phương pháp này có tốc độ huấn luyện nhanh và phù hợp với các tập dữ liệu lớn \cite{mikolov2013efficient}.

\item \textbf{Skip-Gram}: Mô hình thực hiện quá trình ngược lại, sử dụng từ trung tâm để dự đoán các từ ngữ cảnh xung quanh. Skip-Gram thường cho kết quả tốt hơn đối với các từ xuất hiện ít và có khả năng học biểu diễn ngữ nghĩa chính xác hơn CBOW, tuy nhiên chi phí tính toán cao hơn \cite{mikolov2013efficient}.
\end{itemize}
\begin{figure}[htbp]
	\centering
	\includegraphics[width=0.8\textwidth]{2.1}
	\caption{Mô hình CBOW và Skip-Gram.}
	\label{fig:bieudo5}
\end{figure}

Một ưu điểm quan trọng của Word2Vec là khả năng biểu diễn mối quan hệ ngữ nghĩa giữa các từ thông qua các phép toán trên vectơ. Ví dụ, các từ mang ý nghĩa gần nhau như happy, glad và joyful sẽ có các vectơ nằm gần nhau trong không gian đặc trưng. Điều này giúp các mô hình học máy và học sâu khai thác được thông tin ngữ nghĩa mà các phương pháp thống kê truyền thống như TF-IDF khó có thể biểu diễn đầy đủ \cite{mikolov2013efficient}.

Trong phạm vi đề tài, Word2Vec đóng vai trò là một phương pháp biểu diễn văn bản sử dụng trong các phương pháp học sâu CNN và LSTM.

\section{Các mô hình học máy}

\subsection{Logistic Regression}
\subsubsection{Logistic Regression}

Hồi quy Logistic là một thuật toán học máy phân loại mang tính thống kê, được sử dụng để dự đoán xác suất của một biến phụ thuộc phụ thuộc vào các biến độc lập đầu vào \cite{hastie2009elements}. Mặc dù bản chất ban đầu dùng cho bài toán phân loại nhị phân, thuật toán này có thể mở rộng hiệu quả cho bài toán phân loại đa nhãn (Multi-label Classification) thông qua chiến lược Một-đối-Tất cả (One-vs-Rest) \cite{hastie2009elements}

\subsubsection{Hàm Sigmoid và cơ chế kích hoạt xác suất}

Để chuyển đổi giá trị đầu ra tuyến tính thành một giá trị xác suất nằm trong khoảng (0, 1), thuật toán sử dụng hàm Sigmoid được định nghĩa qua công thức:
\begin{equation}
	g(z) = \frac{1}{1 + e^{-z}}
\end{equation}

Trong đó, z là tổ hợp tuyến tính của các đặc trưng đầu vào x và hệ thống trọng số w cùng sai số b:
\begin{equation}
	z = w^T x + b = w_1 x_1 + w_2 x_2 + \dots + w_n x_n + b
\end{equation}

Các thành phần trong phương trình:
\begin{itemize}
	\item $x = [x_1, x_2, \dots, x_n]^T$: Vector chứa các đặc trưng đầu vào (features).
	\item $w = [w_1, w_2, \dots, w_n]^T$: Vector chứa các trọng số (weights) thể hiện mức độ đóng góp của từng đặc trưng.
	\item $b$: Thành phần sai số hệ thống (bias).
\end{itemize}

Đặc điểm hình học và toán học của hàm Sigmoid:
\begin{itemize}
	\item Khi $z \to \infty \implies e^{-z} \to 0 \implies g(z) \to 1$
	
	\item Khi $z \to -\infty \implies e^{-z} \to +\infty \implies g(z) \to 0$
	
	\item Khi $z = 0 \implies e^0 = 1 \implies g(z) = 0.5$
\end{itemize}

Nhờ đường cong mượt mà có hình chữ S, hàm Sigmoid nén toàn bộ không gian số thực vào miền giá trị (0,1). Giá trị đầu ra này được diễn giải trực tiếp thành xác suất điều kiện để mẫu dữ liệu đầu vào thuộc về lớp tích cực (lớp 1):

$P(y = 1|x) = g(z)$

$P(y = 0|x) = 1 - g(z)$

\textbf{Hàm mất mát (Binary Cross-Entropy Loss):} Để tối ưu hóa các trọng số trong quá trình huấn luyện bài toán phân loại đa nhãn, hàm mất mát Entropi chéo nhị phân được áp dụng độc lập cho từng nhãn lớp $j$ \cite{hastie2009elements}:

$$
\mathcal{L}({w}) = -\frac{1}{N} \sum_{i=1}^{N} [y_i \log(\hat{y}_i) + (1 - y_i) \log(1 - \hat{y}_i)]
$$

Trong đó, $N$ là tổng số mẫu huấn luyện, $y_i \in \{0, 1\}$ là nhãn thực tế và $\hat{y}_i = \sigma(z_i)$ là xác suất dự đoán của mô hình \cite{hastie2009elements}.

\subsubsection{Sự tương quan mềm dẻo}

Khái niệm "sự tương quan mềm dẻo" trong Hồi quy Logistic thể hiện qua cách mô hình thiết lập ranh giới phân lớp dựa trên lý thuyết xác suất, thay vì đưa ra các lát cắt cứng nhắc (hard thresholding) như một số thuật toán phân loại dạng cây quyết định hoặc thuật toán Perceptron truyền thống cổ điển \cite{hastie2009elements}.
\begin{itemize}
\item \textbf{Biểu diễn Log-Odds tuyến tính}: Hồi quy Logistic giả định rằng logarit của tỷ số khả năng (log-odds hay giá trị logit) có mối tương quan tuyến tính biến thiên liên tục với các đặc trưng độc lập \cite{hastie2009elements}. Công thức log-odds được xác định qua biểu thức:
$$
\ln \left( \frac{P(y = 1 \mid \mathbf{x})}{1 - P(y = 1 \mid \mathbf{x})} \right) = \mathbf{w}^T \mathbf{x} + b
$$

Mối tương quan này giúp mô hình giữ được bản chất toán học tuyến tính đơn giản ở không gian log-odds, nhưng lại tạo ra một ranh giới chuyển đổi cực kỳ mượt mà, uyển chuyển ở không gian xác suất thực tế \cite{hastie2009elements}.

\item \textbf{Ranh giới quyết định mềm (Soft Decision Boundary)}: Thay vì chỉ khẳng định tuyệt đối một mẫu dữ liệu thuộc lớp nào, mô hình cung cấp mức độ tự tin thông qua giá trị xác suất cụ thể \cite{hastie2009elements}. Tại vùng lân cận ranh giới quyết định nơi xác suất tiến gần về vị trí cân bằng (0.5), giá trị xác suất thay đổi tiệm tiến, giúp phản ánh chính xác sự mơ hồ hoặc dữ liệu nhiễu tại vùng biên của ngôn ngữ mạng xã hội \cite{hastie2009elements}.
\end{itemize}

Sự mềm dẻo này cho phép người vận hành linh hoạt thay đổi ngưỡng phân loại (classification threshold - mặc định là 0.5 để tối ưu hóa mô hình theo các mục tiêu nhiệm vụ khác nhau, như ưu tiên giảm thiểu sai sót loại 1 (False Positive) hoặc sai sót loại 2 (False Negative) \cite{hastie2009elements}

\subsection{Random Forest}

\subsubsection{Học tổ hợp (Ensemble Learning)}

Học tổ hợp là một kỹ thuật nâng cao trong học máy, thực hiện kết hợp dự đoán từ nhiều mô hình thành phần (base models hoặc weak learners) để tạo ra một mô hình tổng hợp duy nhất có hiệu suất tối ưu hơn \cite{hastie2009elements}. Mục tiêu cốt lõi của học tổ hợp là tận dụng "trí tuệ tập thể" nhằm giảm thiểu sai số do phương sai lớn (overfitting/quá khớp) hoặc độ chệch lớn (underfitting/chưa khớp), từ đó nâng cao độ chính xác và khả năng khái quát hóa của mô hình trên dữ liệu thực tế \cite{hastie2009elements}.

\subsubsection{Cơ chế bỏ phiếu số đông}

Bỏ phiếu số đông là một cơ chế tích hợp kết quả đầu ra được sử dụng phổ biến trong các bài toán phân loại (classification) của học tổ hợp.

Cách thức hoạt động: Trong cấu trúc này, mỗi mô hình dự đoán thành phần sẽ hoạt động độc lập và đưa ra một lựa chọn nhãn lớp (class label) cụ thể cho mẫu dữ liệu đầu vào. Hệ thống sẽ tiến hành đếm số lượng phiếu bầu cho từng nhãn lớp. Nhãn lớp nào nhận được số lượng phiếu bầu lớn nhất (chiếm đa số) từ các mô hình thành phần sẽ được chọn làm kết quả dự đoán cuối cùng của toàn bộ hệ thống tổ hợp.

Ví dụ: Giả sử một hệ thống gồm 3 mô hình độc lập thực hiện phân loại một hình ảnh. Mô hình 1 và Mô hình 2 dự đoán kết quả là lớp "Mèo", mô hình 3 dự đoán kết quả là lớp "Chó". Theo cơ chế Hard Voting, lớp "Mèo" nhận được 2/3 số phiếu và được chọn làm kết quả cuối cùng. Cơ chế này khác với Soft Voting (kết hợp dựa trên việc tính trung bình xác suất dự đoán của các lớp).

\subsubsection{Decision Tree}

Khi xây dựng cây quyết định, thuật toán sẽ lặp qua tất cả các đặc trưng và tính toán chỉ số độ lợi thông tin (Information Gain) cho mỗi đặc trưng, sau đó chọn các đặc trưng có Information Gain cao nhất để phân chia dữ liệu tại nút (node) hiện tại \cite{hastie2009elements}.

Information Gain đo lường sự giảm Entropy (độ hỗn loạn/độ không chính xác) của dữ liệu sau khi được phân chia bởi một đặc trưng cụ thể, thể hiện qua công thức:
$$
\text{Gain}(D, A) = \text{Entropy}(D) - \sum_{v \in \text{Values}(A)} \frac{|D_v|}{|D|} \text{Entropy}(D_v)
$$

Entropy càng thấp tức là dữ liệu càng được phân loại một cách chính xác và có độ thuần túy (purity) cao \cite{hastie2009elements}.

\subsubsection{Random Forest}

Random Forest là một thuật toán học máy linh hoạt khi được sử dụng cho các bài toán phân loại, dự đoán giá trị hồi quy tuyến tính và các tác vụ đa đầu ra. Ý tưởng của RF là sử dụng một tập hợp Decision Tree (Cây quyết định), mỗi phân loại được huấn luyện trên các phần khác nhau của tập dữ liệu. Sau đó RF sẽ lấy kết quả phân loại dựa trên các cây phân loại.
\begin{figure}[htbp]
	\centering
	\includegraphics[width=0.8\textwidth]{2.2.png}
	\caption{Mô hình Random Forest.}
	\label{fig:bieudo6}
\end{figure}

Thay vì sử dụng tất cả các đặc trưng để xây dựng cây quyết định, Random Forest sẽ chọn một số lượng ngẫu nhiên các đặc trưng để xem xét tại mỗi node.

Cơ chế hoạt động chi tiết của Random Forest:
\begin{itemize}
\item \textbf{Tạo mẫu ngẫu nhiên (Bootstrapping)}: Từ bộ dữ liệu huấn luyện gốc, tạo ra nhiều tập con dữ liệu bằng cách lấy mẫu ngẫu nhiên với phép lặp (có thể lấy lại một phần tử đã chọn) \cite{hastie2009elements}. Khi tập mẫu được rút ra từ một tập huấn luyện của một cây với sự thay thế (bagging), thì theo ước tính có khoảng 1/3 các phần tử không nằm trong mẫu này \cite{hastie2009elements}. Điều này có nghĩa là chỉ có khoảng 2/3 các phần tử trong tập huấn luyện tham gia vào trong các tính toán của chúng ta, và 1/3 các phần tử này được gọi là dữ liệu ngoài bao (out-of-bag - OOB) \cite{hastie2009elements}. Dữ liệu out-of-bag được sử dụng để ước lượng lỗi tạo ra từ việc kết hợp các kết quả từ các cây tổng hợp trong Random Forest cũng như dùng để ước tính độ quan trọng thuộc tính (variable importance) \cite{hastie2009elements}.

\item \textbf{Xây dựng cây quyết định (Decision Tree)}:
\begin{itemize}
	\item [$\circ$] Chọn Ngẫu Nhiên Các Đặc Trưng (Feature Selection): Tại mỗi nút của cây, chọn ngẫu nhiên một tập hợp con nhỏ các đặc trưng để tìm ra đặc trưng tốt nhất để chia nhánh. Điều này giúp giảm tương quan giữa các cây quyết định [1].
	\item [$\circ$] Chia Dữ Liệu: Sử dụng tiêu chí như Gini impurity hoặc thông tin entropy để tính toán Information Gain và tìm ra cách chia tốt nhất dựa trên các đặc trưng đã chọn \cite{hastie2009elements}.
	\item [$\circ$] Quá trình lặp: Lặp lại quá trình tạo mẫu ngẫu nhiên và xây dựng cây quyết định cho đến khi đạt được số lượng cây quyết định mong muốn trong rừng \cite{hastie2009elements}.
\end{itemize}
\item \textbf{Voting hoặc Averaging}: Khi thực hiện dự đoán, Random Forest sử dụng phương pháp voting (bỏ phiếu - áp dụng cơ chế Hard Voting hoặc Soft Voting trong trường hợp phân loại) hoặc averaging (tính trung bình trong trường hợp hồi quy) từ các cây thành viên \cite{hastie2009elements}. Điều này giúp tạo ra một mô hình tổng hợp mạnh mẽ và giảm phương sai (variance) \cite{hastie2009elements}.
\end{itemize}

\subsection{Support Vector Machine (Linear SVC)}

\subsubsection{Tổng quan về Support Vector Machine (Linear SVC)}

Support Vector Machine (SVM) là một thuật toán học máy có giám sát mạnh mẽ, được ứng dụng rộng rãi cho cả bài toán phân loại và hồi quy. Trong đó, LinearSVC (Linear Support Vector Classification) là một biến thể cụ thể của SVM, được thiết kế tối ưu để giải quyết các bài toán phân loại tuyến tính. Mục tiêu cốt lõi của LinearSVC không chỉ đơn thuần là tìm ra một ranh giới phân chia các lớp dữ liệu, mà là tìm ra một ranh giới "tốt nhất" - mang lại khả năng tổng quát hóa cao nhất trên các dữ liệu chưa từng quan sát.

\subsubsection{Khái niệm siêu phẳng (Hyperplane)}

Trong không gian học máy, ranh giới phân chia các lớp dữ liệu được gọi là một Siêu phẳng (Hyperplane) \cite{hastie2009elements}. Hình dạng và số chiều của siêu phẳng phụ thuộc trực tiếp vào số lượng đặc trưng (features) của tập dữ liệu \cite{hastie2009elements}: Trong không gian 2 chiều (dữ liệu có 2 đặc trưng), siêu phẳng là một đường thẳng \cite{hastie2009elements}. Trong không gian 3 chiều (dữ liệu có 3 đặc trưng), siêu phẳng là một mặt phẳng \cite{hastie2009elements}. Trong không gian n chiều (với n > 3), siêu phẳng là một không gian con có n-1 chiều \cite{hastie2009elements}. Về mặt toán học, phương trình tổng quát của một siêu phẳng phân tách tuyến tính được định nghĩa là \cite{hastie2009elements}:
$$
w^T x + b = 0
$$

Trong đó:
\begin{itemize}
\item $w$: Vectơ pháp tuyến (trọng số) xác định hướng của siêu phẳng.

\item $x$: Vector đặc trưng của điểm dữ liệu đầu vào.

\item $b$: Sai số (bias) xác định vị trí của siêu phẳng so với gốc tọa độ.
\end{itemize}

Mô hình sẽ phân loại một điểm dữ liệu mới bằng cách đưa giá trị $x$ vào phương trình \cite{hastie2009elements}. Nếu $w^T x + b > 0$, điểm dữ liệu thuộc về một lớp (ví dụ: +1); ngược lại, nếu $w^T x + b < 0$, nó thuộc về lớp kia (ví dụ: -1) \cite{hastie2009elements}.

\subsubsection{Tối đa hóa khoảng cách (Max-Margin) và Support Vectors}

Điều làm nên sự khác biệt và sức mạnh của SVM chính là cơ chế Tối đa hóa khoảng cách (Max-Margin) \cite{hastie2009elements}. Trong một tập dữ liệu phân tách tuyến tính được, có thể tồn tại vô số siêu phẳng đi qua giữa hai lớp \cite{hastie2009elements}. Tuy nhiên, SVM chọn siêu phẳng sao cho khoảng cách (margin) từ siêu phẳng đó đến các điểm dữ liệu gần nhất của cả hai lớp là lớn nhất \cite{hastie2009elements}. Support Vectors (Vectơ hỗ trợ): Là các điểm dữ liệu nằm gần siêu phẳng phân tách nhất \cite{hastie2009elements}. Đây là những điểm cực kỳ quan trọng vì nếu loại bỏ hoặc di chuyển các điểm này, vị trí của siêu phẳng sẽ thay đổi \cite{hastie2009elements}. Các điểm dữ liệu nằm xa ranh giới không có bất kỳ ảnh hưởng nào đến việc xây dựng mô hình \cite{hastie2009elements}. Margin (Biên phân cực): Là khoảng cách giữa hai siêu phẳng biên song song với siêu phẳng ranh giới và đi qua các Support Vectors \cite{hastie2009elements}. Phương trình của hai siêu phẳng biên này thường được chuẩn hóa thành \cite{hastie2009elements}:
$$ w^T x + b = 1 $$
$$ w^T x + b = -1 $$

Khoảng cách hình học giữa hai biên này được tính toán bằng công thức:
$$
Margin = \frac{2}{||w||}
$$

Để đạt được khoảng cách (margin) cực đại, thuật toán SVM cần phải tối thiểu hóa độ chuẩn (norm) của vectơ trọng số ||w||, đồng thời phải thỏa mãn điều kiện ràng buộc toán học là tất cả các điểm dữ liệu phải nằm đúng phía của biên phân loại \cite{hastie2009elements}:
$$
y_i (w^T x_i + b) \ge 1 \quad \forall i
$$

Việc tối đa hóa lề (max-margin) giúp mô hình LinearSVC xây dựng được một ranh giới quyết định an toàn nhất, tạo ra khoảng đệm đủ lớn để hấp thụ các nhiễu loạn trong dữ liệu thực tế, từ đó giảm thiểu đáng kể nguy cơ quá khớp (overfitting) và duy trì sự ổn định của hệ thống \cite{hastie2009elements}.

\section{Các mô hình học sâu}

\subsection{CNN}

Mạng thần kinh tích chập (CNN - Convolution Neural Network) là một loại mô hình học sâu, cụ thể là một loại thuật toán học sâu rất phù hợp để phân tích dữ liệu trực quan. CNN - đôi khi được gọi là mạng convnet - sử dụng các nguyên tắc từ đại số tuyến tính, đặc biệt là các phép toán tích chập, để trích xuất các đặc điểm và xác định các mẫu trong hình ảnh. Mặc dù CNN chủ yếu được sử dụng để xử lý hình ảnh nhưng chúng cũng có thể được điều chỉnh để hoạt động với âm thanh và dữ liệu tín hiệu khác.

\subsubsection{Cấu trúc của CNN}

\begin{figure}[H]
	\centering
	\includegraphics[width=0.8\textwidth]{2.3}
	\caption{Kiến trúc của CNN.}
	\label{fig:bieudo7}
\end{figure}

Một CNN thường bao gồm một số lớp, có thể được phân loại thành ba nhóm: lớp tích chập, lớp gộp và lớp được kết nối đầy đủ. Khi dữ liệu đi qua các lớp này, độ phức tạp của CNN tăng lên, cho phép CNN liên tục xác định các phần lớn hơn của hình ảnh và nhiều tính năng trừu tượng hơn.

\textbf{Convolution Layer (Lớp tích chập)}
\begin{itemize}
	\item Lớp tích chập là khối xây dựng cơ bản của CNN và là nơi diễn ra phần lớn các phép tính. Lớp này sử dụng bộ lọc hoặc hạt nhân - một ma trận trọng số nhỏ - để di chuyển qua trường tiếp nhận của hình ảnh đầu vào nhằm phát hiện sự hiện diện của các tính năng cụ thể.
	\item Quá trình bắt đầu bằng cách trượt hạt nhân theo chiều rộng và chiều cao của hình ảnh, cuối cùng quét toàn bộ hình ảnh qua nhiều lần lặp. Tại mỗi vị trí, một tích số chấm được tính giữa trọng số của hạt nhân và giá trị pixel của hình ảnh bên dưới hạt nhân. Điều này biến đổi hình ảnh đầu vào thành một tập hợp các bản đồ đặc trưng hoặc các đặc điểm phức tạp, mỗi đặc điểm thể hiện sự hiện diện và cường độ của một đặc điểm nhất định tại các điểm khác nhau trong hình ảnh.
	\item CNN thường bao gồm nhiều lớp chập xếp chồng lên nhau. Thông qua kiến trúc phân lớp này, CNN diễn giải dần dần thông tin hình ảnh có trong dữ liệu hình ảnh thô. Trong các lớp trước, CNN xác định các đặc điểm cơ bản như cạnh, kết cấu hoặc màu sắc. Các lớp sâu hơn nhận đầu vào từ bản đồ đặc trưng của các lớp trước, cho phép chúng phát hiện các mẫu, vật thể và cảnh phức tạp hơn.
\end{itemize}

\textbf{Pooling Layer (Lớp gộp)}
\begin{itemize}
	\item Lớp gộp của CNN là một thành phần quan trọng theo sau lớp tích chập. Tương tự như lớp tích chập, các hoạt động của lớp gộp liên quan đến quá trình quét trên hình ảnh đầu vào, nhưng chức năng của nó thì khác.
	\item Lớp tổng hợp nhằm mục đích giảm tính chiều của dữ liệu đầu vào trong khi vẫn giữ lại thông tin quan trọng, do đó cải thiện hiệu quả tổng thể của mạng. Điều này thường đạt được thông qua việc lấy mẫu xuống: giảm số lượng điểm dữ liệu trong đầu vào.
	\item Đối với CNN, điều này thường có nghĩa là giảm số lượng pixel được sử dụng để thể hiện hình ảnh. Hình thức gộp phổ biến nhất là gộp tối đa, giữ lại giá trị tối đa trong một cửa sổ nhất định (tức là kích thước hạt nhân) trong khi loại bỏ các giá trị khác. Một kỹ thuật phổ biến khác, được gọi là gộp trung bình, có cách tiếp cận tương tự nhưng sử dụng giá trị trung bình thay vì giá trị tối đa.
	\item Việc lấy mẫu xuống làm giảm đáng kể tổng số tham số và tính toán. Ngoài việc nâng cao hiệu quả, điều này còn tăng cường khả năng khái quát hóa của mô hình. Các mô hình ít phức tạp hơn với các tính năng cấp cao hơn thường ít bị trang bị quá mức – một hiện tượng xảy ra khi một mô hình nhận biết nhiễu và các chi tiết quá cụ thể trong dữ liệu huấn luyện của nó, ảnh hưởng tiêu cực đến khả năng khái quát hóa thông tin mới, chưa được nhìn thấy.
	\item Việc giảm kích thước không gian của việc biểu diễn có một nhược điểm tiềm ẩn, đó là mất một số thông tin. Tuy nhiên, chỉ học những đặc điểm nổi bật nhất của dữ liệu đầu vào thường là đủ cho các tác vụ như phát hiện đối tượng và phân loại hình ảnh. Fully connected layer.
\end{itemize}

\textbf{Fully Connected Layer (Lớp kết nối đầy đủ)}
\begin{itemize}
	\item Lớp được kết nối đầy đủ đóng một vai trò quan trọng trong các giai đoạn cuối cùng của CNN, nơi nó chịu trách nhiệm phân loại hình ảnh dựa trên các đặc điểm được trích xuất ở các lớp trước đó. Thuật ngữ kết nối đầy đủ có nghĩa là mỗi nơ-ron trong một lớp được kết nối với mỗi nơ-ron ở lớp tiếp theo, cho phép CNN xem xét đồng thời tất cả các tính năng khi đưa ra quyết định phân loại cuối cùng.
	\item Không phải tất cả các lớp trong CNN đều được kết nối đầy đủ. Bởi vì các lớp được kết nối đầy đủ có nhiều tham số, nên việc áp dụng phương pháp này trên toàn bộ mạng sẽ tạo ra mật độ không cần thiết, làm tăng nguy cơ trang bị quá mức và khiến mạng rất tốn kém khi đào tạo về bộ nhớ và tính toán. Việc giới hạn số lượng các lớp được kết nối đầy đủ sẽ cân bằng hiệu quả tính toán và khả năng khái quát hóa với khả năng tìm hiểu các mẫu phức tạp.
\end{itemize}

\subsection{LSTM}

Mạng bộ nhớ dài - ngắn (LSTM - Long Short-Term Memory) là một biến thể kiến trúc bậc cao của Mạng nơ-ron hồi quy (RNN). LSTM được thiết kế đặc biệt để khắc phục triệt để vấn đề "triệt tiêu đạo hàm" (vanishing gradient) của RNN truyền thống, mang lại khả năng ghi nhớ và chọn lọc các thông tin phụ thuộc xa trong chuỗi dữ liệu thời gian. Trong lĩnh vực xử lý ngôn ngữ tự nhiên, đặc biệt là nhận diện cảm xúc (thông qua các nghiên cứu thực tiễn như WASSA-2018 và SemEval-2018), LSTM cho thấy sức mạnh vượt trội nhờ khả năng nắm bắt được ngữ cảnh và trật tự từ vựng của cả một câu dài.

Kiến trúc LSTM: Thay vì chỉ có một mạng nơ-ron đơn giản, mỗi ô nhớ (Cell) của LSTM được cấu tạo phức tạp bởi một trạng thái ô (tương đương với luồng bộ nhớ dài hạn) và ba cổng kiểm soát thông tin: cổng quên, cổng đầu vào và cổng đầu ra.
\begin{itemize}
	\item Trạng thái ô (Cell State) : lưu trữ trạng thái của một chuỗi, do đó nó có khả năng giữ hoặc quên một số thông tin nhất định.
	\item Cổng đầu vào (Input Gate): quyết định thông tin mới nào được lưu trữ trong ô.
	\item Cổng đầu ra (Output Gate): xác định trạng thái ẩn tiếp theo.
	\item Cổng quên (Forget Gate): quyết định thông tin nào nên bỏ đi hoặc giữ lại.
\end{itemize}

\textbf{Hoạt động của LSTM}

Giả sử chúng ta có một chuỗi các từ $(w_1, w_2, w_3, \dots, w_n)$ và chúng ta đang xử lý từng từ một trong chuỗi đó. Biểu thị trạng thái của LSTM tại bước thời gian t là $(h_t, C_t)$, trong đó $h_t$ là trạng thái ẩn và $C_t$ là trạng thái ô.

Bước 1: LSTM nhận vectơ đầu vào ($x_t$) và trạng thái trước đó ($h_{t-1}, C_{t-1}$).

Bước 2: Cổng quên ($f_t$) quyết định thông tin nào cần loại bỏ khỏi trạng thái ô. Nó sử dụng vectơ đầu vào và trạng thái ẩn trước đó để tạo một số từ 0 đến 1 cho mỗi số ở trạng thái ô $C_{t-1}$. Giá trị gần 0 nghĩa là thông tin đó sẽ bị loại bỏ, còn giá trị gần 1 nghĩa là thông tin đó sẽ được giữ lại.
\begin{equation}
	f_t = \sigma(W_f \cdot [h_{t-1}, x_t] + b_f)
\end{equation}

Bước 3: Cổng đầu vào ($i_t$) quyết định thông tin mới nào sẽ được lưu trữ ở trạng thái ô. Nó có hai phần: Lớp sigmoid được gọi là "lớp cổng đầu vào" quyết định giá trị nào sẽ cập nhật
\begin{equation}
	i_t = \sigma(W_i [h_{t-1}, x_t] + b_i)
\end{equation}

Lớp tanh tạo ra một vectơ gồm các trạng thái ô mới ($\tilde{C}_t$) có thể được thêm vào $C_t$
\begin{equation}
	\tilde{C}_t = \tanh(W_c \cdot [h_{t-1}, x_t] + b_c)
\end{equation}

Bước 4: Cập nhật trạng thái ô cũ ($C_{t-1}$) sang trạng thái ô mới ($C_t$). Trạng thái ô cũ được nhân với $f_t$ để quên đi những điều chúng ta đã quyết định quên trước đó. Sau đó, chúng ta thêm các trạng thái ô mới, được điều chỉnh bởi $i_t$ để tạo ra $C_t$
\begin{equation}
	C_t = f_t \cdot C_{t-1} + i_t \cdot \tilde{C}_t
\end{equation}

Bước 5: Quyết định đầu ra. Đầu tiên, chúng ta chạy một lớp sigmoid để quyết định phần nào của trạng thái ô mà chúng ta sẽ xuất ra. Sau đó, chúng ta đặt trạng thái ô thông qua hàm tanh (để đặt các giá trị nằm trong khoảng từ -1 đến 1) và nhân nó với đầu ra $o_t$, chỉ xuất những phần mà chúng ta quyết định giữ lại.
\begin{equation}
	h_t = o_t \cdot \tanh(C_t)
\end{equation}

\begin{figure}[H]
	\centering
	\includegraphics[width=0.8\textwidth]{2.4.png}
	\caption{Mô hình LSTM trong xử lý ngôn ngữ tự nhiên.}
	\label{fig:bieudo8}
\end{figure}

LSTM bao gồm bốn phần chính: Đầu vào nhúng từ, mạng di động LSTM, fully connected layer và softmax. Với đầu vào, mỗi ô trong mạng LSTM nhận được một vectơ từ được biểu thị bằng các từ nhúng có dạng $[1\times{n}]$ trong đó $n$ là độ dài cố định của câu. Sau đó, các ô tính toán các giá trị và nhận kết quả dưới dạng vectơ trong mạng di động LSTM. Các vectơ này sẽ được kết nối đầy đủ và các giá trị đầu ra sau đó sẽ chuyển qua hàm softmax để đưa ra phân loại phù hợp cho từng nhãn.

\section{Các chỉ số đánh giá}

Trong bài toán phân loại đa nhãn cảm xúc với 28 trạng thái tâm lý khác nhau, tập dữ liệu thực tế thường xảy ra tình trạng mất cân bằng nghiêm trọng (ví dụ: nhãn "neutral" hoặc "approval" xuất hiện dày đặc, trong khi "grief" hoặc "pride" lại rất hiếm). Do đó, chỉ số Độ chính xác tổng thể (Accuracy) truyền thống không thể phản ánh đúng hiệu năng thực sự của hệ thống. Để có cái nhìn khách quan và toàn diện nhất, đề tài sử dụng tập hợp các chỉ số đánh giá chuyên sâu bao gồm: Precision, Recall và đặc biệt là F1-Score thông qua hai cơ chế tính toán Micro và Macro.

Trước khi định nghĩa các chỉ số, ta xét các giá trị thống kê cơ bản trong Ma trận nhầm lẫn (Confusion Matrix) đối với một nhãn cụ thể:
\begin{itemize}
	\item TP (True Positive): Số trường hợp mô hình dự đoán đúng nhãn cảm xúc đó.
	
	\item FP (False Positive): Số trường hợp mô hình dự đoán có nhãn cảm xúc đó, nhưng thực tế không phải (Báo động giả).
	\item FN (False Negative): Số trường hợp nhãn cảm xúc đó thực sự tồn tại, nhưng mô hình lại bỏ sót.
\end{itemize}

\subsection{Precision (Độ chính xác)}

Tỷ lệ số dự đoán đúng trên tổng số các dự đoán dương tính của mô hình. Trong nhận diện cảm xúc, Precision cao đồng nghĩa với việc khi hệ thống đã gán một nhãn cảm xúc nào đó (ví dụ: \textit{"vui vẻ"}), thì khả năng cao văn bản đó thực sự mang cảm xúc \textit{"vui vẻ"}.
$$
Precision = \frac{TP}{TP+FP}
$$

\subsection{Recall (Độ bao phủ)}

Tỷ lệ số dự đoán đúng trên tổng số các mẫu dương tính thực tế. Chỉ số này đánh giá khả năng "không bỏ sót" của mô hình. Recall cao cho thấy hệ thống có khả năng nhận diện và gom được hầu hết các bình luận có chứa cảm xúc đang xét.
$$
Recall = \frac{TP}{TP+FN}
$$

\subsection{Đánh giá tổng thể hệ thống}

Do hệ thống phải dự đoán cùng lúc 28 nhãn cảm xúc, F1-Score của từng nhãn cần được tổng hợp lại thành một thước đo duy nhất để so sánh hiệu năng giữa các thuật toán (như Logistic Regression, LinearSVC, CNN, LSTM) và để tinh chỉnh siêu tham số. Đề tài sử dụng hai phương pháp tính trung bình:
\begin{itemize}
	\item Micro F1-Score: Tính toán tổng số lượng TP, FP, FN trên toàn bộ 28 nhãn rồi mới áp dụng công thức F1. Chỉ số này chịu ảnh hưởng lớn bởi các nhãn đa số. Nếu hệ thống dự đoán tốt các nhãn xuất hiện nhiều, Micro F1 sẽ cao.
	
	\item Macro F1-Score: Tính toán chỉ số F1 hoàn toàn độc lập cho từng nhãn cảm xúc, sau đó lấy trung bình cộng của 28 giá trị F1 đó. Phương pháp này coi trọng tất cả các nhãn ngang nhau, bất kể số lượng mẫu của nhãn đó nhiều hay ít.
\end{itemize}

Trong khuôn khổ của bài toán, Macro F1-Score được lựa chọn làm thước đo tối ưu hóa trọng tâm. Mục tiêu tối thượng của hệ thống không phải là học vẹt các cảm xúc phổ biến, mà là phải có khả năng nhận diện công bằng và nhạy bén với cả những trạng thái tâm lý thiểu số tinh vi nhất.

\subsection{Đường cong học tập (Learning Curve) và cơ chế dừng sớm}

Khác với các thuật toán học máy truyền thống, các mô hình Học sâu (như 1D-CNN và LSTM) tiến hành tối ưu hóa trọng số liên tục qua hàng loạt chu kỳ lặp (Epoch). Do đó, Đường cong học tập (Learning Curve) được sử dụng làm công cụ chuẩn đoán trực quan cốt lõi để giám sát động thái học tập của mạng nơ-ron.

Đường cong học tập biểu diễn sự biến thiên của hai dải thông số chính trên cả tập huấn luyện (Training set) và tập kiểm định (Validation set) qua từng chu kỳ:
\begin{itemize}
	\item \textbf{Hàm mất mát (Loss)}: Đối với bài toán phân loại đa nhãn, hàm mất mát nhị phân (Binary Cross-Entropy) thường được theo dõi. Đồ thị lý tưởng sẽ có xu hướng giảm dần, phản ánh sự suy giảm sai số dự đoán của mô hình.
	
	\item \textbf{Chỉ số đánh giá (Metric)}: Các chỉ số hiệu năng (như Macro F1-Score) được ghi nhận song song, với kỳ vọng đồ thị sẽ tăng tịnh tiến và duy trì ở trạng thái bão hòa.
\end{itemize}

Thông qua việc quan sát hình thái của đường cong, sự tương tác giữa dữ liệu và thuật toán được thể hiện qua ba trạng thái cơ bản:
\begin{itemize}
	\item \textbf{Chưa khớp (Underfitting)}: Đường cong mất mát của cả hai tập dữ liệu đều neo ở mức cao và không có xu hướng giảm. Điều này cho thấy mạng nơ-ron chưa đủ độ sâu hoặc độ phức tạp để bóc tách được các đặc trưng ngữ nghĩa của văn bản.
	
	\item \textbf{Hội tụ lý tưởng (Good Fit)}: Hai đường cong huấn luyện và kiểm định cùng giảm dần và bám sát nhau, tiệm cận về một ngưỡng sai số cực tiểu và ổn định.
	
	\item \textbf{Quá khớp (Overfitting)}: Đây là rủi ro đặc biệt nghiêm trọng khi xử lý ngôn ngữ tự nhiên trên các tập dữ liệu mất cân bằng. Tại trạng thái này, đường mất mát trên tập huấn luyện tiếp tục lao dốc (thậm chí tiệm cận 0), nhưng đường mất mát trên tập kiểm định lại chững lại, đảo chiều và tăng vọt. Hiện tượng này khẳng định hệ thống đang "học vẹt" cả những nhiễu loạn vô nghĩa của tập huấn luyện, dẫn đến sự sụp đổ về khả năng khái quát hóa khi gặp dữ liệu thực tế.
\end{itemize}

Nhằm triệt tiêu rủi ro quá khớp và tối ưu hóa thời gian tính toán, đề tài áp dụng cơ chế \textbf{Dừng sớm (Early Stopping)} kết hợp cùng đường cong học tập. Thay vì ép buộc thuật toán chạy hết một số lượng Epoch cố định định trước, hệ thống sẽ liên tục giám sát đồ thị Validation Loss. Nếu phát hiện hàm mất mát trên tập kiểm định không tiếp tục giảm sau một số chu kỳ chịu đựng (Patience), quá trình huấn luyện sẽ bị ngắt tự động. Thuật toán sẽ quay lui (restore) và trích xuất lại bộ trọng số tại "điểm ngọt" (Sweet spot) – nơi Validation Loss đạt mức thấp nhất, đảm bảo mô hình mang ra kiểm thử là phiên bản mạnh mẽ và khách quan nhất. 
% Chương 3

\chapter{THỰC NGHIỆM} % Tên của chương

\label{Chapter3} % Để trích dẫn chương này ở chỗ nào đó trong bài, hãy sử dụng lệnh \ref{Chapter3} 

%----------------------------------------------------------------------------------------

\section{Phân tích dữ liệu}

\subsection{Khái quát về tập dữ liệu GoEmotions}

Đề tài này tiến hành thực nghiệm trên tập dữ liệu GoEmotions, một trong những bộ dữ liệu gán nhãn thủ công lớn nhất dành cho bài toán nhận diện đa nhãn cảm xúc, được giới thiệu và phát triển bởi các nhà nghiên cứu tại Google \cite{demszky2020goemotions}. Bộ dữ liệu bao gồm tổng cộng 58.009 dòng bình luận tiếng Anh được thu thập từ nền tảng Reddit. Điểm đặc biệt làm nên giá trị của GoEmotions chính là hệ thống phân loại cảm xúc vô cùng chi tiết khi toàn bộ các bình luận được gán nhãn thủ công cho 27 danh mục cảm xúc khác nhau cùng với một nhãn Trung tính, tổng cộng là 28 nhãn \cite{demszky2020goemotions}. Hệ thống phân loại này bao phủ đa dạng các sắc thái tâm lý từ tích cực, tiêu cực cho đến những trạng thái cảm xúc không rõ ràng, đáp ứng hoàn hảo mục tiêu nghiên cứu của đề tài.

Về quy trình xây dựng, dữ liệu thô ban đầu được trích xuất từ các cộng đồng phổ biến trên Reddit trong giai đoạn từ năm 2005 đến tháng 1 năm 2019 \cite{demszky2020goemotions}. Trên thực tế, nền tảng mạng xã hội Reddit vốn nổi tiếng với sự thiên lệch lớn về mặt nhân khẩu học, khi đối tượng đóng góp nội dung chủ yếu nghiêng về người dùng nam trẻ tuổi, dẫn đến việc không phản ánh được sự đa dạng của toàn cầu (Duggan và Smith, 2013) \cite{demszky2020goemotions}. Đồng thời, cấu trúc vận hành của nền tảng này cũng dễ tạo ra sự lệch lạc nghiêm trọng về ngôn từ độc hại và xúc phạm (Mohan và cộng sự, 2017) \cite{demszky2020goemotions}. Do đặc thù đó, nội dung trên Reddit thường được sử dụng để nghiên cứu về chứng trầm cảm (Pirina và Cöltekin, 2018) \cite{demszky2020goemotions}, các hành vi vi phạm nhỏ (Breitfeller và cộng sự, 2019) \cite{demszky2020goemotions}, hay thậm chí nghiên cứu của Yanardag và Rahwan (2018) đã chỉ ra tác động tiêu cực của việc sử dụng dữ liệu Reddit thiên lệch bằng cách huấn luyện thành công một chatbot mang hành vi của "kẻ tâm thần" \cite{demszky2020goemotions}.

Chính vì vậy, để giải quyết những lo ngại này và cho phép xây dựng các mô hình phân loại cảm xúc mang tính đại diện rộng rãi, các tác giả của GoEmotions đã thực hiện một loạt các biện pháp sàng lọc, tiền xử lý dữ liệu nghiêm ngặt nhằm đảm bảo tập dữ liệu cuối cùng không củng cố các định kiến ngôn ngữ chung hoặc cụ thể về mặt cảm xúc \cite{demszky2020goemotions}.
\begin{itemize}
	\item \textbf{Giảm thiểu ngôn từ độc hại}: Loại bỏ các subreddit không an toàn (NSFW) và các cộng đồng có trên 10\% bình luận chứa token xúc phạm, người lớn. Các từ chửi thề thông thường được giữ lại vì chúng cốt lõi để mô hình học các cảm xúc tiêu cực.
	
	\item \textbf{Kiểm duyệt thủ công}: Rà soát và loại bỏ các bình luận mang tính xúc phạm hướng đến chủng tộc, giới tính, xu hướng tình dục hoặc người khuyết tật.
	
	\item \textbf{Lọc độ dài}: Giữ lại bình luận từ 3 đến 30 token. Tiến hành lấy mẫu giảm với giới hạn chặn dựa trên số lượng bình luận có số token bằng giá trị trung vị.
	
	\item \textbf{Cân bằng tính phân cực và cảm xúc}: Sử dụng một mô hình thử nghiệm huấn luyện trên 2.200 mẫu để ước lượng cảm xúc. Từ đó, loại bỏ các subreddit chứa hơn 30\% bình luận trung tính hoặc ít hơn 20\% bình luận tích cực/tiêu cực/mơ hồ, đồng thời thực hiện downsampling để giảm định kiến cảm xúc và cân bằng tỷ lệ giữa các subreddit.
	
	\item \textbf{Che giấu thông tin}: Từ 315.000 bình luận sạch sau bộ lọc (thuộc 482 subreddit), các tác giả lấy mẫu ngẫu nhiên để gán nhãn. Tên riêng chỉ người và các thuật ngữ tôn giáo được nhận diện bằng bộ gán nhãn thực thể BERT-based để thay thế bằng các mã token [NAME] và [RELIGION] nhằm triệt tiêu các định kiến phân loại.
\end{itemize}
\begin{figure}[H]
	\centering
	\includegraphics[width=0.8\textwidth]{3.1.png}
	\caption{Minh họa các mẫu văn bản tiếng Anh thô trích xuất từ dữ liệu.}
	\label{fig:bieudo9}
\end{figure}

Sau giai đoạn làm sạch, quá trình gán nhãn được thực hiện bởi đội ngũ gồm 82 người bản xứ nói tiếng Anh đến từ Ấn Độ nhằm đảm bảo tính chính xác tối đa về mặt ngữ nghĩa \cite{demszky2020goemotions}. Quy trình này tuân thủ cơ chế kiểm định nghiêm ngặt: mỗi bình luận ban đầu sẽ được đánh giá độc lập bởi 3 người; nếu không đạt được sự đồng thuận về ít nhất một nhãn, hệ thống sẽ tự động bổ sung thêm 2 người gán nhãn khác. Người gán nhãn được phép chọn nhiều cảm xúc cùng lúc nếu văn bản thể hiện sự đa chiều, hoặc lựa chọn nhãn Neutral trong trường hợp không chắc chắn. Kết quả thống kê cho thấy chất lượng gán nhãn đạt độ đồng thuận rất cao khi có tới 94\% số mẫu (tương đương 54.263 bình luận) đạt ít nhất 2 sự đồng thuận và 31\% số mẫu đạt từ 3 sự đồng thuận trở lên \cite{demszky2020goemotions}. Tuy nhiên, mức độ nhận diện giữa các lớp có sự phân hóa rõ rệt: các trạng thái hiển nhiên như biết ơn (gratitude), ngưỡng mộ (admiration), và thích thú (amusement) đạt độ đồng thuận rất cao; trái lại, các cảm xúc mang tính ngầm ẩn như đau buồn (grief) hay lo lắng (nervousness) lại có mức độ đồng thuận thấp nhất do đòi hỏi nhiều ngữ cảnh phức tạp hơn \cite{demszky2020goemotions}.

\begin{figure}[H]
	\centering
	\includegraphics[width=0.8\textwidth]{3.2.png}
	\caption{Mô tả bộ dữ liệu GoEmotions.}
	\label{fig:bieudo9}
\end{figure}

\begin{figure}[H]
	\centering
	\includegraphics[width=0.8\textwidth]{3.1.1.3.png}
	\caption{Thống kê tổng quan về quy mô và chất lượng gán nhãn của tập dữ liệu GoEmotions.}
	\label{fig:bieudo3.1.1.3}
\end{figure}

\subsection{Phân bố nhãn cảm xúc}

Phân tích thống kê tần suất xuất hiện của 28 nhãn trên tập dữ liệu GoEmotions bộc lộ một thách thức cốt lõi của đề tài: Sự mất cân bằng dữ liệu (Class Imbalance) cực kỳ nghiêm trọng. Cấu trúc phân bố này tuân theo dạng đuôi dài (long-tail distribution), phản ánh đúng đặc tính tự nhiên của ngôn ngữ giao tiếp trên mạng xã hội.

Ở nhóm chiếm ưu thế tuyệt đối, nhãn \textit{Neutral} (Trung tính) có tần suất xuất hiện cao nhất, chiếm tỷ trọng áp đảo trong không gian nhãn. Theo sau đó là các trạng thái cảm xúc mang tính tương tác xã hội tích cực và dễ bộc lộ trên internet như \textit{Approval} (Tán thành), \textit{Admiration} (Ngưỡng mộ), và \textit{Joy} (Vui vẻ). Sự hội tụ tại nhóm nhãn này hoàn toàn hợp lý với văn hóa diễn đàn Reddit, nơi phần lớn bình luận là những câu phản hồi ngắn gọn thể hiện sự đồng tình hoặc khen ngợi đối với bài đăng gốc.

Ở chiều ngược lại, phần "đuôi dài" của phổ phân bố tập trung những cảm xúc có tần suất xuất hiện rất thấp, cá biệt một số nhãn chỉ chiếm chưa tới 1\% tổng số lượng mẫu. Các trạng thái như \textit{Grief} (Đau buồn), \textit{Relief} (Nhẹ nhõm), hay \textit{Pride} (Tự hào) cực kỳ khan hiếm. Sự thiếu hụt này xuất phát từ bản chất tâm lý học: đây là những dòng cảm xúc tinh vi, đòi hỏi những bối cảnh rất cụ thể, mang tính cá nhân cao và hiếm khi được người dùng biểu đạt trực tiếp qua các thông điệp văn bản ngắn trên không gian mạng mở.

Đặc biệt, do tính chất của bài toán phân loại đa nhãn (Multi-label), một bình luận có thể được gán nhiều hơn một cảm xúc (ví dụ: vừa \textit{Amusement} vừa \textit{Approval}). Sự phân bố bất đối xứng và đan xen này đòi hỏi hệ thống thực nghiệm không thể dùng Accuracy làm thước đo chuẩn, mà bắt buộc phải áp dụng các kỹ thuật bù đắp trọng số (Class Weighting) và sử dụng chỉ số Macro F1-Score để đánh giá công bằng năng lực nhận diện của mô hình trên các nhãn thiểu số.
\begin{figure}[H]
	\centering
	\includegraphics[width=0.8\textwidth]{3.3.png}
	\caption{Phân bố nhãn cảm xúc.}
	\label{fig:bieudo10}
\end{figure}

Đáng chú ý, biểu đồ (Hình 3.4) không chỉ phản ánh sự mất cân bằng nghiêm trọng về mặt số lượng (dạng đuôi dài), mà thang màu sắc còn chỉ ra một thách thức kép về mặt ngữ nghĩa. Các nhãn có tần suất xuất hiện cao ở nhóm đầu (như \textit{admiration}, \textit{gratitude}, \textit{amusement}) đều sở hữu dải màu tím đậm, chứng tỏ mức độ đồng thuận đánh giá rất cao. Ngược lại, các nhãn thiểu số nằm ở phần đuôi (như \textit{grief}, \textit{relief}, \textit{pride}) lại mang màu xanh nhạt hoặc trắng, tương ứng với hệ số tương quan rất thấp. Điều này khẳng định rằng: các cảm xúc hiếm gặp trong bộ dữ liệu cũng chính là những trạng thái tâm lý mơ hồ và khó nhận diện nhất đối với cả con người. Đây là cơ sở thực tiễn để đề tài áp dụng kỹ thuật tìm ngưỡng tối ưu tự động (Threshold Tuning) nhằm cứu vãn độ bao phủ (Recall) cho các nhãn thiểu số này trong giai đoạn thực nghiệm mô hình.

\subsection{Tương quan giữa các nhãn}

Bên cạnh việc khảo sát tần suất độc lập, phân tích ma trận tương quan (Correlation Heatmap) kết hợp phân cụm phân cấp (Hierarchical Clustering) dựa trên sự đồng xuất hiện của các nhãn cung cấp một góc nhìn sâu sắc về cấu trúc tâm lý ẩn trong bộ dữ liệu (Hình 3.1.3.1). Kết quả trực quan hóa cho thấy các nhãn cảm xúc không tồn tại biệt lập mà có sự liên kết mạng lưới chặt chẽ, tạo thành các cụm ngữ nghĩa mang tính phân cấp.

Thứ nhất, tồn tại sự tương quan dương (Positive Correlation) mạnh mẽ giữa các cụm cảm xúc có sự tương đồng về sắc thái, thể hiện qua các vùng màu đỏ/hồng trên ma trận. Biểu hiện rõ nhất là sự đồng xuất hiện thường xuyên giữa cụm \textit{Anger} (Tức giận) đi kèm với \textit{Annoyance} (Khó chịu), hoặc cụm \textit{Joy} (Vui vẻ) và \textit{Amusement} (Thích thú). Điều này phản ánh logic tâm lý tự nhiên của con người, khi một trạng thái bức xúc thường kéo theo thái độ gắt gỏng trong ngôn từ. Việc tồn tại các cụm tương quan này giúp các thuật toán học máy (đặc biệt là mạng nơ-ron sâu) có thể chia sẻ trọng số và dễ dàng khái quát hóa các vùng không gian vectơ đặc trưng có ý nghĩa lân cận nhau.

Thứ hai, ma trận cũng chỉ ra sự tương quan âm (Negative Correlation) rõ rệt giữa các nhóm cảm xúc đối lập, thể hiện qua các vùng giao cắt màu xanh dương. Các cặp nhãn như \textit{Sadness} (Buồn bã) và \textit{Joy} (Vui vẻ) gần như loại trừ lẫn nhau và hiếm khi xuất hiện trên cùng một không gian văn bản. Sự phân tách này tạo ra các ranh giới siêu phẳng (hyperplanes) tương đối sắc nét, hỗ trợ các thuật toán học máy tuyến tính (như LinearSVC hay Logistic Regression) bóc tách không gian và hoạt động ổn định hơn.

Thứ ba, sơ đồ nhánh cây (Dendrogram) ở viền biểu đồ đã tự động phân rã 27 nhãn thành 3 cụm lớn. Cấu trúc này khớp hoàn toàn với 3 dải phân cực cảm xúc (Sentiment) cốt lõi bao gồm: Tích cực (màu xanh lá), Tiêu cực (màu xanh dương) và Mơ hồ/Đa nghĩa (màu vàng).

Tóm lại, việc thấu hiểu cấu trúc tương quan đa nhãn không chỉ giải mã được hành vi biểu đạt cảm xúc của tập dữ liệu mà còn là cơ sở khoa học quan trọng. Nó định hướng trực tiếp cho đề tài sử dụng chiến lược One-Vs-Rest kết hợp cùng các hàm kích hoạt Sigmoid độc lập ở lớp đầu ra, cho phép mô hình nhận diện thành công sự đan xen phức tạp của các dòng trạng thái tâm lý

\begin{figure}[H]
	\centering
	\includegraphics[width=0.8\textwidth]{3.1.3.png}
	\caption{Biểu đồ nhiệt (heatmap) thể hiện sự tương quan giữa các đánh giá cho từng cảm xúc. Biểu đồ biểu thị sự phân cụm phân cấp của các điểm đánh giá. Việc gán nhãn cảm xúc được thực hiện từ trước và cho thấy rằng các cụm tương ứng chặt chẽ với các nhóm cảm xúc \cite{demszky2020goemotions}.}
	\label{fig:bieudo9}
\end{figure}

\section{Tiền xử lý}

\subsection{Làm sạch văn bản}

Sau khi thu thập và sàng lọc, các chuỗi văn bản thô từ Reddit vẫn chứa rất nhiều nhiễu ngôn ngữ, từ viết tắt, ký tự đặc biệt và lỗi mã hóa ký tự. Để chuẩn hóa dữ liệu đầu vào cho các mô hình, đề tài đã xây dựng một hàm tiền xử lý chuyên sâu mang tên advanced\_clean\_text. Quy trình này được thực hiện tuần tự qua các bước cốt lõi sau:
\begin{itemize}
	\item \textbf{Chuẩn hóa bảng mã và đồng bộ hóa ký tự:} Nhằm khắc phục triệt để hiện tượng lỗi hiển thị phông chữ (lỗi encoding thường gặp khi xử lý dữ liệu trích xuất từ Reddit), hệ thống sử dụng biểu thức chính quy để đồng bộ hóa mọi dạng ký tự dấu nháy đơn phi chuẩn \verb|[’‘´]| về dạng dấu nháy đơn chuẩn mực \verb|[']|. Đồng thời, thư viện \texttt{unicodedata} với chuẩn chuẩn hóa NFKD được áp dụng kết hợp với mã hóa ASCII để bóc tách, loại bỏ các dấu thanh, dấu phụ và ép toàn bộ chuỗi ký tự về định dạng văn bản ASCII thuần túy mà không làm mất đi ký tự gốc. 
	
	\item \textbf{Chuyển đổi chữ thường:} Toàn bộ chuỗi văn bản được đưa về dạng chữ viết thường. Bước này giúp mô hình đối xử thống nhất các từ có cách viết hoa khác nhau (như "\textit{Awesome}", "\textit{AWESOME}" và "\textit{awesome}"), từ đó làm giảm đáng kể kích thước không gian từ vựng.
	
	\item \textbf{Khai triển từ viết tắt:} Đặc trưng ngôn ngữ mạng xã hội của Reddit là mật độ từ viết tắt rất cao. Hàm tiền xử lý tiến hành tra cứu và giải nén các cụm từ viết tắt phổ biến (ví dụ: \textit{i'm} thành \textit{i am}, \textit{don't} thành \textit{do not}, \textit{it's} thành \textit{it is}) và sử dụng các bộ lọc đuôi cấu trúc (\textit{n't}, \textit{'re}, \textit{'ve}, \textit{'ll}...) để vớt các từ viết tắt còn sót lại, đưa văn bản về đúng cấu trúc ngữ pháp gốc.
	
	\item \textbf{Kỹ thuật nối từ phủ định:} Đây là một kỹ thuật nâng cao được cấy vào hệ thống nhằm bảo toàn ngữ nghĩa biểu cảm. Biểu thức chính quy sẽ tự động nhận diện cấu trúc từ phủ định đi kèm một động từ/tính từ liền sau và nối chúng lại bằng dấu gạch dưới (ví dụ: \textit{not like} thành \textit{not\_like}, \textit{not good} thành \textit{not\_good}). Kỹ thuật này giúp các mô hình học máy nhận biết được các cụm từ mang ý nghĩa đảo ngược, tránh việc đánh giá sai lệch trọng số của các từ tích cực/tiêu cực đơn lẻ.
	
	\item \textbf{Lọc nhiễu ký tự và ép từ lặp vô nghĩa:} Hệ thống sử dụng biểu thức chính quy để loại bỏ toàn bộ các ký tự đặc biệt, chỉ giữ lại khoảng trắng, ký tự chữ cái \verb|[a-z]|, chữ số \verb|[0-9]| và hai dấu biểu cảm cốt lõi là ! và ?. Đồng thời, các ký tự lặp lại vô nghĩa do thói quen gõ chữ của người dùng (ví dụ: \textit{soooo} thành \textit{so}) được ép về dạng gốc nhằm loại bỏ phân tán dữ liệu.
	
	\item \textbf{Chuẩn hóa gốc từ:} Ép toàn bộ các biến thể động từ mang các thì khác nhau (như \textit{loved}, \textit{loving}, \textit{loves}) quay trở về dạng nguyên thể gốc (\textit{love}), giúp tối ưu hóa không gian đặc trưng đầu vào.
\end{itemize}

\begin{figure}[H]
	\centering
	\includegraphics[width=0.8\textwidth]{3.4.png}
	\caption{Văn bản sau khi xử lý.}
	\label{fig:bieudo11}
\end{figure}

\subsection{Biểu diễn văn bản và trích xuất đặc trưng}

\subsubsection{Biểu diễn văn bản bằng phương pháp TF - IDF}

Mục đích của TF-IDF: các thuật toán học máy truyền thống như Logictis Regulation, Support Vector Machine (SVM) hay Random Forest, không thể trực tiếp đọc hiểu các chuỗi văn bản thô. TF-IDF được sử dụng để chuyển đổi các đoạn văn bản thành các vector số học. Phương pháp này giúp mô hình đánh giá được mức độ quan trọng của từng từ vựng đối với một văn bản cụ thể trong một tập ngữ liệu lớn.

Hệ thống tiến hành sử dụng công cụ TfidfVectorizer từ thư viện scikit-learn của Python để trích xuất đặc trưng văn bản. Quá trình cấu hình được tối ưu hóa thông qua các tham số kỹ thuật cụ thể:
\begin{itemize}
	\item \textbf{Cấu hình vùng từ (ngram\_range=(1, 2))}: Cho phép mô hình trích xuất cả từ đơn (Unigram) lẫn cụm hai từ liền kề (Bigram). Việc này giúp hệ thống bắt được ngữ cảnh của các cụm từ quan trọng, đặc biệt là các từ phủ định đã được liên kết ở bước tiền xử lý.
	
	\item \textbf{Giới hạn từ vựng (max\_features=5000)}: Hệ thống chỉ giữ lại 5.000 từ và cụm từ có tần suất xuất hiện tối ưu nhất trên toàn bộ tập dữ liệu Reddit nhằm giảm chiều dữ liệu, loại bỏ nhiễu và tăng tốc độ huấn luyện cho mô hình.
\end{itemize}

\subsubsection{Phương pháp nhúng từ Word2Vec}

Khác với TF-IDF vốn chỉ tạo ra các ma trận thưa (sparse matrix) và bỏ qua mối liên hệ ngữ cảnh giữa các từ, phương pháp Word2Vec được áp dụng cho các cấu trúc mạng nơ-ron học sâu (CNN, LSTM) nhằm chuyển đổi văn bản thành các ma trận dày đặc (dense matrix). Việc này giúp hệ thống mô hình hóa được sự tương đồng về mặt ý nghĩa (semantic similarity) giữa các từ vựng.

Quá trình trích xuất đặc trưng được cấu hình với các tham số kỹ thuật cụ thể như sau:
\begin{itemize}
	\item \textbf{Kiến trúc huấn luyện}: Sử dụng mô hình Skip-Gram (sg=1). Kiến trúc này được chọn vì hiệu quả vượt trội trong việc học biểu diễn các từ hiếm (infrequent words), vốn xuất hiện rất nhiều trong dữ liệu bình luận ngắn.
	
	\item \textbf{Số chiều không gian (Vector size)}: d = 100. Mỗi từ được ánh xạ thành một vectơ toán học chứa 100 đặc trưng số thực. Kích thước này được tính toán để cân bằng giữa khả năng biểu diễn ngữ nghĩa và chi phí tính toán phần cứng.
	
	\item \textbf{Kích thước cửa sổ trượt (Window size)}: window = 5. Mô hình sẽ đánh giá mối quan hệ của một từ trung tâm với 5 từ liền trước và 5 từ liền sau nó để học ngữ cảnh.
	
	\item \textbf{Ngưỡng tần suất (Min count)}: min\_count = 2. Các từ vựng (ví dụ như lỗi gõ phím ngẫu nhiên hoặc tên riêng) chỉ xuất hiện duy nhất 1 lần trong toàn bộ 58.000 câu bình luận sẽ bị loại bỏ hoàn toàn khỏi không gian từ vựng để tránh nhiễu mô hình.
\end{itemize}

Kết quả đầu ra của quá trình này là một ma trận nhúng (Embedding Matrix), đóng vai trò như một lớp từ điển số học để nạp trực tiếp vào Layer đầu tiên (Embedding Layer) của các mô hình học sâu ở phần 3.3.

\section{Huấn luyện mô hình}

Toàn bộ quá trình lập trình, cài đặt thực nghiệm và huấn luyện các mô hình học máy được nhóm tiến hành trên nền tảng điện toán đám mây kết hợp với phần cứng máy tính cá nhân. Môi trường thực thi cụ thể được thiết lập như sau:
\begin{itemize}
\item \textbf{Phần cứng máy tính cá nhân (PC)}: Máy tính để bàn chạy hệ điều hành Windows 10 Pro 64-bit, trang bị bộ vi xử lý Intel® Pentium® CPU G4400 @ 3.30GHz và bộ nhớ RAM 8GB.

\item \textbf{Môi trường thực thi đám mây}: Google Colaboratory (Colab) được sử dụng làm môi trường thực thi chính, khai thác sức mạnh của bộ xử lý GPU nhằm tăng tốc độ huấn luyện cho các mô hình học sâu (Deep Learning).

\item \textbf{Ngôn ngữ và Nền tảng}: Python 3, kết hợp với các thư viện học máy cốt lõi bao gồm Scikit-learn, TensorFlow và Keras.
\end{itemize}

\subsection{Mô hình hồi quy Logisic}

Đầu vào của hệ thống tiếp tục là không gian đặc trưng văn bản thưa được số hóa qua phương pháp thống kê TF-IDF. Mặc dù Hồi quy Logistic nguyên bản được thiết kế cho bài toán phân loại nhị phân, hệ thống đã khéo léo ứng dụng chiến lược One-Vs-Rest (OvR) để mở rộng năng lực thuật toán, xây dựng 28 mô hình phân lớp tuyến tính độc lập nhằm giải quyết trọn vẹn bài toán đa nhãn (tương ứng với 28 sắc thái cảm xúc).

Điểm nhấn cốt lõi của quy trình huấn luyện này là cơ chế thực nghiệm dò tìm siêu tham số C (Nghịch đảo của hệ số điều chuẩn - Inverse of Regularization Strength). Hệ thống sẽ tự động duyệt qua các mức điều chuẩn từ nặng đến nhẹ (ví dụ: 0.01, 0.1, 1, 10) để tìm ra trạng thái cân bằng lý tưởng nhất: vừa đủ khắt khe để triệt tiêu hiện tượng quá khớp (Overfitting) trên không gian ma trận TF-IDF khổng lồ, vừa đủ linh hoạt để mô hình học được các ranh giới từ vựng quan trọng. Bằng cách tính toán và so sánh trực tiếp độ đo Macro F1-Score trên tập kiểm định (Validation set), thuật toán sẽ chốt lại cấu hình trọng số ưu việt nhất và gán làm mô hình chính thức, tạo tiền đề vững chắc cho giai đoạn hiệu chỉnh ngưỡng quyết định (Threshold Tuning) phía sau.

Để khắc phục tình trạng các nhãn cảm xúc thiểu số bị bỏ qua hoàn toàn khi sử dụng ngưỡng quyết định mặc định (0.5), hệ thống áp dụng kỹ thuật dò tìm ngưỡng tối ưu tự động trên tập kiểm định (Validation Set). Cụ thể, sau khi mô hình sinh ra ma trận xác suất thô, thuật toán sẽ quét độc lập qua 28 nhãn và thử nghiệm liên tục các giá trị ngưỡng từ 0.1 đến 1.0. Tại mỗi nhãn, mức ngưỡng mang lại điểm F1-Score cao nhất sẽ được tự động trích xuất và lưu lại thành một từ điển cấu hình. Bộ 28 ngưỡng tối ưu cục bộ này sau đó được áp dụng trực tiếp lên tập kiểm thử mù (Test Set) để nhị phân hóa kết quả dự đoán. Kỹ thuật này giúp hệ thống linh hoạt nới lỏng ranh giới (hạ ngưỡng) cho các cảm xúc khó đoán và siết chặt (tăng ngưỡng) với các nhãn dễ đoán, từ đó tối đa hóa độ bao phủ (Recall) trên tập dữ liệu mất cân bằng nghiêm trọng.

\subsection{Mô hình Random Forest}

Đầu vào của hệ thống là các đoạn văn bản đã được số hóa thành ma trận đặc trưng toán học thông qua phương pháp thống kê TF-IDF. Để giải quyết bài toán phân loại đa nhãn, hệ thống áp dụng chiến lược One-Vs-Rest (OvR) để khởi tạo và huấn luyện song song 28 mô hình "Rừng ngẫu nhiên" (Random Forest) độc lập. Mỗi mô hình lõi chịu trách nhiệm nhận diện một cảm xúc riêng biệt và được tinh chỉnh chặt chẽ với 100 cây quyết định, đồng thời giới hạn độ sâu (max\_depth) cùng số lượng mẫu tối thiểu tại các ngọn lá để triệt tiêu hiện tượng học vẹt (Overfitting) trên dữ liệu mất cân bằng. Ở giai đoạn dự đoán, hệ thống sẽ tổng hợp kết quả phân loại độc lập từ toàn bộ 28 khu rừng, từ đó đưa ra quyết định cuối cùng và cho phép gán đồng thời nhiều sắc thái tâm lý đan xen cho cùng một bình luận. 

Để khắc phục tình trạng các nhãn cảm xúc thiểu số bị bỏ qua hoàn toàn khi sử dụng ngưỡng quyết định mặc định (0.5), hệ thống áp dụng kỹ thuật dò tìm ngưỡng tối ưu tự động trên tập kiểm định (Validation Set). Cụ thể, sau khi mô hình sinh ra ma trận xác suất thô, thuật toán sẽ quét độc lập qua 28 nhãn và thử nghiệm liên tục các giá trị ngưỡng từ -1 đến 0.95. Tại mỗi nhãn, mức ngưỡng mang lại điểm F1-Score cao nhất sẽ được tự động trích xuất và lưu lại thành một từ điển cấu hình. Bộ 28 ngưỡng tối ưu cục bộ này sau đó được áp dụng trực tiếp lên tập kiểm thử mù (Test Set) để nhị phân hóa kết quả dự đoán. Kỹ thuật này giúp hệ thống linh hoạt nới lỏng ranh giới (hạ ngưỡng) cho các cảm xúc khó đoán và siết chặt (tăng ngưỡng) với các nhãn dễ đoán, từ đó tối đa hóa độ bao phủ (Recall) trên tập dữ liệu mất cân bằng nghiêm trọng.

\subsection{Mô hình Support Vector Machine (SVM)}

Đầu vào của hệ thống là ma trận đặc trưng thưa (sparse matrix) được số hóa thông qua phương pháp thống kê TF-IDF. Để giải quyết trọn vẹn bài toán phân loại đa nhãn, hệ thống áp dụng chiến lược One-Vs-Rest (OvR) bao bọc bên ngoài thuật toán lõi LinearSVC, qua đó thiết lập 28 siêu mặt phẳng phân cách (Hyperplanes) độc lập tương ứng với 28 sắc thái cảm xúc.

Đặc biệt, quá trình huấn luyện được tích hợp cơ chế dò tìm siêu tham số tự động để tìm ra giá trị điều chuẩn C (Soft Margin Penalty) tối ưu nhất. Thuật toán sẽ lặp qua một tập hợp các giá trị C khác nhau để tinh chỉnh "độ mềm" của lề phân loại — cân bằng giữa việc mở rộng ranh giới để khái quát hóa dữ liệu mạng xã hội và việc siết chặt ranh giới để bám sát tập huấn luyện. Điểm số Macro F1-Score được tính toán trên tập kiểm định (Validation set) sẽ đóng vai trò làm thước đo định lượng. Cấu hình C mang lại giá trị Macro F1 cao nhất sẽ được hệ thống tự động lựa chọn và gán làm mô hình chính thức, sẵn sàng cho bước tinh chỉnh ngưỡng quyết định (Threshold Tuning) tiếp theo.

Để khắc phục tình trạng các nhãn cảm xúc thiểu số bị bỏ qua hoàn toàn khi sử dụng ngưỡng quyết định mặc định (0.5), hệ thống áp dụng kỹ thuật dò tìm ngưỡng tối ưu tự động trên tập kiểm định (Validation Set). Cụ thể, sau khi mô hình sinh ra ma trận xác suất thô, thuật toán sẽ quét độc lập qua 28 nhãn và thử nghiệm liên tục các giá trị ngưỡng từ 0.1 đến 1.0. Tại mỗi nhãn, mức ngưỡng mang lại điểm F1-Score cao nhất sẽ được tự động trích xuất và lưu lại thành một từ điển cấu hình. Bộ 28 ngưỡng tối ưu cục bộ này sau đó được áp dụng trực tiếp lên tập kiểm thử mù (Test Set) để nhị phân hóa kết quả dự đoán. Kỹ thuật này giúp hệ thống linh hoạt nới lỏng ranh giới (hạ ngưỡng) cho các cảm xúc khó đoán và siết chặt (tăng ngưỡng) với các nhãn dễ đoán, từ đó tối đa hóa độ bao phủ (Recall) trên tập dữ liệu mất cân bằng nghiêm trọng.

\subsection{Mô hình CNN (1D - CNN)}

Đầu vào của hệ thống là chuỗi các từ vựng đã được ánh xạ thành các vectơ số thực thông qua mạng nhúng từ (Word Embeddings). Hệ thống triển khai mạng thần kinh tích chập một chiều (1D-CNN) với nhiều bộ lọc có kích thước khác nhau quét dọc theo câu nhằm bóc tách các đặc trưng cụm từ (n-gram) cục bộ trong bình luận tiếng Anh. Các tín hiệu ngữ nghĩa nổi bật nhất sau đó được chắt lọc qua lớp Gộp tối đa (Max Pooling) và ghép nối thành một vectơ đặc trưng tổng hợp duy nhất đại diện cho toàn bộ câu. Đặc biệt, để phù hợp với bài toán phân loại đa nhãn, lớp nơ-ron đầu ra (Output Layer) được thiết kế kết hợp cùng hàm kích hoạt Sigmoid nhằm tính toán xác suất xuất hiện độc lập của từng cảm xúc, từ đó cho phép một bình luận có thể được gán đồng thời nhiều sắc thái tâm lý khác nhau.

Để khắc phục tình trạng các nhãn cảm xúc thiểu số bị bỏ qua hoàn toàn khi sử dụng ngưỡng quyết định mặc định (0.5), hệ thống áp dụng kỹ thuật dò tìm ngưỡng tối ưu tự động trên tập kiểm định (Validation Set). Cụ thể, sau khi mô hình sinh ra ma trận xác suất thô, thuật toán sẽ quét độc lập qua 28 nhãn và thử nghiệm liên tục các giá trị ngưỡng từ 0.1 đến 1.0. Tại mỗi nhãn, mức ngưỡng mang lại điểm F1-Score cao nhất sẽ được tự động trích xuất và lưu lại thành một từ điển cấu hình. Bộ 28 ngưỡng tối ưu cục bộ này sau đó được áp dụng trực tiếp lên tập kiểm thử mù (Test Set) để nhị phân hóa kết quả dự đoán. Kỹ thuật này giúp hệ thống linh hoạt nới lỏng ranh giới (hạ ngưỡng) cho các cảm xúc khó đoán và siết chặt (tăng ngưỡng) với các nhãn dễ đoán, từ đó tối đa hóa độ bao phủ (Recall) trên tập dữ liệu mất cân bằng nghiêm trọng.

\subsection{Mô hình LSTM}

Đầu vào của hệ thống là chuỗi các từ vựng đã được ánh xạ thành các vectơ số thực thông qua mạng nhúng từ (Word Embeddings). Chuỗi vectơ này sẽ được truyền lần lượt qua các ô nhớ LSTM theo trật tự thời gian (kiến trúc Many-to-One). Sau khi quét đến từ cuối cùng của câu, trạng thái ẩn cuối cùng mang theo toàn bộ thông tin ngữ cảnh tổng hợp sẽ được truyền qua một lớp kết nối đầy đủ (Fully connected layer). Do tính chất của bài toán phân loại đa nhãn (Multi-label Classification), mạng nơ-ron đầu ra sẽ sử dụng hàm kích hoạt Sigmoid để đánh giá độc lập xác suất xuất hiện của từng nhãn cảm xúc, từ đó cho phép một bình luận có thể được gán đồng thời nhiều sắc thái tâm lý khác nhau. 

Để khắc phục tình trạng các nhãn cảm xúc thiểu số bị bỏ qua hoàn toàn khi sử dụng ngưỡng quyết định mặc định (0.5), hệ thống áp dụng kỹ thuật dò tìm ngưỡng tối ưu tự động trên tập kiểm định (Validation Set). Cụ thể, sau khi mô hình sinh ra ma trận xác suất thô, thuật toán sẽ quét độc lập qua 28 nhãn và thử nghiệm liên tục các giá trị ngưỡng từ 0.1 đến 1.0. Tại mỗi nhãn, mức ngưỡng mang lại điểm F1-Score cao nhất sẽ được tự động trích xuất và lưu lại thành một từ điển cấu hình. Bộ 28 ngưỡng tối ưu cục bộ này sau đó được áp dụng trực tiếp lên tập kiểm thử mù (Test Set) để nhị phân hóa kết quả dự đoán. Kỹ thuật này giúp hệ thống linh hoạt nới lỏng ranh giới (hạ ngưỡng) cho các cảm xúc khó đoán và siết chặt (tăng ngưỡng) với các nhãn dễ đoán, từ đó tối đa hóa độ bao phủ (Recall) trên tập dữ liệu mất cân bằng nghiêm trọng.

\section{Đánh giá mô hình}

Để đo lường và so sánh khách quan hiệu năng giữa các thuật toán học máy truyền thống và các mạng học sâu trong bài toán phân loại đa nhãn với 28 sắc thái cảm xúc, nhóm tiến hành sử dụng các phương pháp và độ đo đánh giá tiêu chuẩn sau:

\subsection{Báo cáo phân loại}

Đây là bảng tổng hợp chi tiết hiệu suất của mô hình, cung cấp ba chỉ số cốt lõi cho từng nhãn cảm xúc bao gồm: Precision (Độ chuẩn xác), Recall (Độ bao phủ) và F1-Score (Trung bình điều hòa của Precision và Recall). Đặc biệt, do bộ dữ liệu GoEmotions bị mất cân bằng phân lớp nghiêm trọng, độ đo Accuracy (Độ chính xác tổng thể) sẽ gây ra những đánh giá sai lệch (ảo tưởng sức mạnh mô hình). Do đó, đề tài quyết định sử dụng Macro F1-Score làm thước đo tối thượng để so sánh các mô hình. Chỉ số này tính toán F1 độc lập cho từng nhãn rồi lấy trung bình cộng, bắt buộc mô hình phải dự đoán tốt cả những nhãn thiểu số (như \textit{grief}, \textit{pride}) thì mới đạt được điểm cao.
 
\subsection{Ma trận nhầm lẫn}

Đây là công cụ trực quan hóa dạng ma trận hai chiều, so sánh trực tiếp giữa nhãn thực tế (True Label) và nhãn do mô hình dự đoán (Predicted Label). Đối với không gian đặc trưng phức tạp của văn bản, ma trận nhầm lẫn không chỉ giúp tính toán số lượng mẫu đúng/sai, mà còn giúp nhóm nhận diện được hiện tượng mô hình thường xuyên nhầm lẫn giữa các cảm xúc có chung vĩ độ tâm lý hoặc sắc thái gần giống nhau (ví dụ: nhầm lẫn giữa \textit{annoyance} - \textit{khó chịu} và \textit{anger} - \textit{giận dữ})

\subsection{Đường cong học tập (Learning Curve / Loss Curve)}

Độ đo này được áp dụng đặc biệt cho quá trình huấn luyện các mạng thần kinh học sâu (CNN, LSTM). Đồ thị biểu diễn sự biến thiên của hàm mất mát (Training Loss / Validation Loss) và điểm số đánh giá qua từng chu kỳ huấn luyện (epoch). Việc theo dõi đường cong này giúp hệ thống đánh giá được tốc độ hội tụ của mô hình, đồng thời phát hiện sớm các hiện tượng học vẹt, quá khớp (Overfitting) hay chưa khớp (Underfitting) để từ đó có những điều chỉnh kịp thời về tham số Dropout hay L2-Regularization.
% Chương 4

\chapter{KẾT QUẢ VÀ THẢO LUẬN} % Tên của chương

\label{Chapter4} % Để trích dẫn chương này ở chỗ nào đó trong bài, hãy sử dụng lệnh \ref{Chapter4}

%----------------------------------------------------------------------------------------

\section{Kết quả}

\subsection{Hồi quy Logistic + TF - IDF}

\begin{figure}[H]
	\centering
	\includegraphics[width=0.8\textwidth]{4.1.png}
	\caption{Classification Report của Hồi quy Logistic + TF - IDF.}
	\label{fig:bieudo12}
\end{figure}

\textbf{Hiệu suất tổng thể}: Mô hình đạt điểm số Macro F1-Score là 0.43. Đây là một mức điểm phản ánh đúng thực tế của bài toán phân loại đa nhãn phức tạp (28 lớp) trên dữ liệu văn bản ngắn, đồng thời cho thấy sự khó khăn của thuật toán tuyến tính khi phải xử lý không gian đặc trưng thưa thớt của mạng xã hội.

\textbf{Điểm sáng ở các nhãn đa số}: Thuật toán hoạt động cực kỳ hiệu quả ở những nhãn cảm xúc có số lượng mẫu dồi dào (support cao). Điển hình như nhãn gratitude (327 mẫu) đạt F1-Score xuất sắc lên tới 0.89, hay nhãn amusement (316 mẫu) đạt 0.78. Điều này chứng tỏ khi được cung cấp đủ dữ liệu, chiến lược One-Vs-Rest của Logistic Regression học được các ranh giới quyết định rất tốt.

\textbf{Hạn chế ở các nhãn thiểu số}: Ngược lại, hiện tượng mất cân bằng dữ liệu đã tác động tiêu cực đến năng lực dự đoán của mô hình ở các lớp hiếm gặp. Nhãn grief (chỉ có 7 mẫu test) có điểm F1-Score chạm đáy ở mức 0.13, hay relief (19 mẫu test) chỉ đạt 0.16. Do tập huấn luyện quá mỏng, mô hình không thể khái quát hóa được các từ khóa cốt lõi của những cảm xúc này, dẫn đến độ bao phủ (recall) rất thấp.
\begin{figure}[H]
	\centering
	\includegraphics[width=0.8\textwidth]{4.2.png}
	\caption{Confusion Matrix của Hồi quy Logistic + TF - IDF.}
	\label{fig:bieudo13}
\end{figure}

Mô hình dự đoán cực kỳ nhạy bén và chuẩn xác ở các nhãn có dữ liệu dồi dào (Amusement, Gratitude, Love). Ngược lại, nó gần như "bị mù" và không thể nhận diện được các nhãn thiểu số (Grief, Pride, Relief) do số ca đoán trúng (True Positive) chỉ đếm trên đầu ngón tay.

Ở nhãn Trung tính (Neutral): Khi đối mặt với các câu bình luận chứa cảm xúc lấp lửng, mơ hồ hoặc chứa từ lóng chưa học được, mô hình có xu hướng an toàn đẩy hết kết quả về nhãn Neutral. Điều này khiến lượng dự đoán nhầm (False Positive) ở nhãn Trung tính vọt lên mức cao nhất toàn hệ thống (1367 mẫu).

\subsection{Random Forest + TF - IDF}

\begin{figure}[H]
	\centering
	\includegraphics[width=0.8\textwidth]{4.3.png}
	\caption{Classification Report của Random Forest + TF - IDF.}
	\label{fig:bieudo14}
\end{figure}
\begin{figure}[H]
	\centering
	\includegraphics[width=0.8\textwidth]{4.4.png}
	\caption{Confusion Matrix của Random Forest + TF - IDF.}
	\label{fig:bieudo15}
\end{figure}

\textbf{Hiệu năng tổng thể}: Mô hình đạt điểm Macro F1-Score là 0.40. Dù điểm số tổng thể có phần khiêm tốn hơn so với các mô hình tuyến tính, nhưng đây vẫn là một kết quả đáng ghi nhận đối với thuật toán Tree-based (dựa trên cây) khi phải xử lý ma trận thưa (sparse matrix) khổng lồ của TF-IDF trên bài toán 28 nhãn mất cân bằng nghiêm trọng.

\textbf{Hạn chế bộc lộ trên các nhãn thiểu số}: Random Forest gặp khó khăn cực lớn với các nhãn hiếm (như grief, pride, relief). Đối chiếu sang Ma trận nhầm lẫn, số ca đoán trúng (True Positive) ở các nhãn này gần như bằng 0. Do số lượng mẫu quá ít, các cây quyết định không thể phân nhánh đủ sâu để tạo ra luật dự đoán, dẫn đến độ bao phủ (Recall) sụp đổ hoàn toàn.

\textbf{Phong độ ổn định ở nhãn đa số}: Bù lại, đối với các cảm xúc có tần suất xuất hiện dày đặc như amusement, joy hay anger, Random Forest phát huy tối đa sức mạnh của học tổ hợp (Ensemble Learning). Việc tổng hợp kết quả bình chọn từ hàng trăm cây độc lập giúp mô hình nhận diện cực kỳ ổn định, bắt trúng các từ khóa cốt lõi và duy trì độ chuẩn xác (Precision) cao.

\textbf{Hiện tượng "đổ đống" vào nhãn Trung tính (Neutral)}: Giống như các mô hình truyền thống khác, giới hạn của phương pháp Bag-of-Words khiến Random Forest không hiểu được ngữ cảnh vĩ mô. Điểm yếu này thể hiện rõ nhất trên Ma trận nhầm lẫn khi lượng dự đoán nhầm (False Positive) ở nhãn Neutral vọt lên mức cao nhất. Khi gặp các câu lấp lửng, thiếu từ khóa rõ ràng, các cây quyết định có xu hướng "bỏ phiếu trắng", đẩy toàn bộ kết quả về nhãn Trung tính để an toàn.

\subsubsection{LinearSVC + TF - IDF}

\begin{figure}[H]
	\centering
	\includegraphics[width=0.8\textwidth]{4.5.png}
	\caption{Classification Report của LinearSVC + TF - IDF.}
	\label{fig:bieudo16}
\end{figure}
\begin{figure}[H]
	\centering
	\includegraphics[width=0.8\textwidth]{4.6.png}
	\caption{Confusion Matrix của LinearSVC + TF - IDF.}
	\label{fig:bieudo17}
\end{figure}

\textbf{Hiệu năng tổng thể}: Mô hình đạt điểm Macro F1-Score là 0.43. Sự phân bổ hiệu suất giữa các nhãn lại cho thấy sự vượt trội rõ rệt về mặt thuật toán của SVM đối với bài toán mất cân bằng dữ liệu.

\textbf{Sức mạnh trên các nhãn thiểu số}: LinearSVC sử dụng cơ chế tối ưu hóa lề (margin) dựa trên các Vector Hỗ trợ (Support Vectors). Nhờ việc tập trung kẻ siêu mặt phẳng phân cách (Hyperplane) cứng dựa trên các điểm dữ liệu nằm sát biên, SVM đã "vớt vát" thành công độ bao phủ (Recall) của các nhãn cực hiếm, đẩy điểm F1 của nhóm này lên mức khả quan hơn đáng kể.

\textbf{Phong độ ổn định ở nhãn đa số}: Đối với các cảm xúc có tần suất xuất hiện dày đặc như amusement, gratitude hay love, LinearSVC vẫn duy trì được độ chuẩn xác (Precision) rất cao. Khả năng xử lý tốt ma trận thưa (sparse matrix) 5000 chiều của TF-IDF giúp thuật toán này bóc tách ranh giới của các từ khóa đặc trưng một cách sắc bén.

\textbf{Hạn chế}: Dù đã cải thiện được việc nhận diện nhãn thiểu số, SVM vẫn vấp phải rào cản chung của các mô hình học máy truyền thống khi không thể hiểu được ngữ cảnh vĩ mô và thứ tự từ vựng trong câu (do bản chất của TF-IDF là Bag-of-Words). Điều này khiến mô hình vẫn bị nhầm lẫn (False Positive) nhiều ở nhãn Neutral khi gặp các câu có cấu trúc lấp lửng.

\subsection{Multi-Kernel 1D CNN + Word2Vec}

\begin{figure}[H]
	\centering
	\includegraphics[width=0.8\textwidth]{4.1.4.4.png}
	\caption{Classification Report của Multi-Kernel 1D CNN + Word2Vec.}
	\label{fig:bieudo4.1.4.4}
\end{figure}

Đánh giá hiệu suất phân loại dựa trên F1-Score Kết quả dự đoán trên tập kiểm thử (Test Set) cung cấp một cái nhìn chi tiết về sức mạnh cũng như giới hạn của mô hình trên tập dữ liệu phân bố lệch (imbalanced data).
\begin{itemize}
	\item Sự cân bằng giữa Precision và Recall: Điểm Macro Precision đạt 0.4 và Macro Recall đạt 0.44 cho thấy thuật toán Dò tìm ngưỡng tối ưu (Threshold Tuning) đã hoạt động chính xác. Mô hình không bị thiên lệch theo hướng dự đoán bừa (khiến Recall cao, Precision thấp) hay quá khắt khe (khiến Precision cao, Recall thấp) mà đã thiết lập được một điểm cắt (cut-off) hài hòa nhất.
	
	\item Những nhãn hiếm gặp trong tập huấn luyện (chỉ xuất hiện vài chục mẫu) như \textit{Grief}, \textit{Relief}, chưa có đủ khối lượng ngữ cảnh để mô hình học được đặc trưng, dẫn đến F1-Score xấp xỉ 0. Khi tính trung bình không trọng số (Macro Average), điểm số của các nhãn thiểu số này đã kéo tụt tổng điểm hệ thống.
	
	\item Các điểm sáng trong nhận diện cảm xúc: Bù lại, với các nhãn có khối lượng dữ liệu lớn hoặc mang các từ khóa ngữ nghĩa rõ ràng, mạng CNN thể hiện sức mạnh nhận diện vượt trội. Điển hình như nhãn \textit{Gratitude} (Biết ơn) đạt điểm F1 lên tới 0.89, \textit{Love} (Yêu thương) đạt 0.78, \textit{Neutral} (Trung lập) đạt 0.65, và \textit{Remorse} (Hối hận) đạt 0.62. Điều này khẳng định rằng đối với những cụm từ mang sắc thái biểu cảm định hình rõ, các bộ lọc Multi-Kernel của mạng CNN hoàn toàn đủ khả năng quét và trích xuất thành công.
\end{itemize}
\begin{figure}[H]
	\centering
	\includegraphics[width=0.8\textwidth]{4.1.4.2.png}
	\caption{Biểu đồ lịch sử huấn luyện của Multi-Kernel 1D CNN + Word2Vec.}
	\label{fig:bieudo16}
\end{figure}

Từ hình 4.1.4.2, đánh giá biểu đồ lịch sử huấn luyện cho thấy mô hình Multi-Kernel 1D CNN có khả năng trích xuất đặc trưng ngôn ngữ và hội tụ rất tốt trong giai đoạn đầu. Từ Epoch 1 đến Epoch 4, hàm mất mát trên tập kiểm định (Validation Loss) giảm đều đặn và chạm mức thấp nhất ở khoảng 0.103. Tương ứng với đó, độ đo Validation AUC tăng trưởng mạnh mẽ và đạt đỉnh tại 84.5\% (0.845). Mức AUC này là một minh chứng rõ ràng cho thấy cấu trúc mạng tích chập đề xuất kết hợp với biểu diễn từ vựng (Word Embeddings) đã phân biệt thành công các dải phổ cảm xúc đa dạng của văn bản.

Tuy nhiên, từ Epoch thứ 5 trở đi, hiện tượng quá khớp (Overfitting) bắt đầu xuất hiện. Trong khi Training Loss tiếp tục suy giảm tuyến tính xuống mức 0.06 và Training AUC tiệm cận mức 0.95 (mô hình học thuộc dữ liệu), thì Validation Loss lại có xu hướng nảy ngược đi lên (đạt >0.12) và Validation AUC đi ngang. Mặc dù vậy, nhờ việc áp dụng thành công cơ chế Dừng sớm (Early Stopping) với tham số restore\_best\_weights, hệ thống đã tự động bảo toàn được bộ trọng số tối ưu nhất tại bước sóng phân kỳ (Epoch 4), đảm bảo mô hình giữ được khả năng tổng quát hóa khi dự đoán trên dữ liệu thực tế.
\begin{figure}[H]
	\centering
	\includegraphics[width=0.8\textwidth]{4.1.4.3.png}
	\caption{Confusion matrix của Multi-Kernel 1D CNN + Word2Vec.}
	\label{fig:bieudo16}
\end{figure}

Nhìn chung, bộ ma trận nhầm lẫn cho thấy mô hình xử lý bài toán phân loại đa nhãn khá hiệu quả, đặc biệt xuất sắc ở các nhóm cảm xúc có từ vựng biểu đạt rõ ràng như \textit{Amusement}, \textit{Gratitude}, \textit{Love} và \textit{Neutral} (tỷ lệ True Positive cao, False Negative thấp). Sự áp đảo của các ô True Negative phản ánh đúng bản chất thưa thớt của dữ liệu văn bản thực tế, nơi mỗi bình luận thường chỉ mang 1 đến 2 nhãn cảm xúc. Tuy nhiên, kết quả cũng chỉ ra hai hạn chế cục bộ: các nhãn thiểu số (\textit{Grief}, \textit{Pride}) gần như không thể nhận diện do hiện tượng thiếu hụt dữ liệu huấn luyện (Data Sparsity), trong khi các cảm xúc có ranh giới ngữ nghĩa mập mờ (\textit{Annoyance}, \textit{Realization}) lại có tỷ lệ nhận diện nhầm (False Positive) tương đối cao. Tổng thể, nhờ sức mạnh biểu diễn từ vựng và chiến lược tinh chỉnh ngưỡng cắt, mạng học sâu đã phân tách tốt các đặc trưng cốt lõi dù phải đối mặt với tập dữ liệu mất cân bằng nghiêm trọng.

\subsection{LSTM + Word2Vec}

\begin{figure}[H]
	\centering
	\includegraphics[width=0.8\textwidth]{4.1.5.1.png}
	\caption{Classification Report của LSTM + Word2Vec.}
	\label{fig:bieudo4.1.5.1}
\end{figure}

Đánh giá hiệu suất phân loại dựa trên F1-Score của mô hình LSTM + Word2Vec: Dựa trên Bảng báo cáo phân loại (Classification Report), kết quả thực nghiệm của mô hình LSTM kết hợp Word2Vec cho thấy hệ thống đạt điểm Macro F1-Score là 0.39 và Micro F1-Score là 0.52.
\begin{itemize}
	\item Điểm sáng ở các nhãn đa số: Mô hình LSTM phát huy sức mạnh ở các nhãn có khối lượng dữ liệu dồi dào và ngữ cảnh rõ ràng, điển hình như nhãn gratitude đạt F1-Score 0.88, amusement đạt 0.80 và love đạt 0.76. Điều này cho thấy khả năng ghi nhớ trật tự từ vựng theo thời gian của LSTM (thông qua kiến trúc Many-to-One) đã giúp mô hình bắt trúng các ngữ cảnh biểu đạt cảm xúc mạnh.
	
	Hạn chế ở các nhãn thiểu số: Tương tự như các mô hình trước, hiện tượng mất cân bằng dữ liệu tiếp tục là rào cản lớn. Các nhãn thiểu số như grief có điểm F1-Score rơi xuống 0.00, pride đạt 0.12 và relief đạt 0.10. Sự thiếu hụt dữ liệu khiến các trạng thái ẩn (hidden states) của LSTM không có đủ bối cảnh để học được các mẫu đặc trưng của những cảm xúc tinh vi này, dẫn đến độ bao phủ (Recall) cực kỳ thấp.
\end{itemize}

\begin{figure}[H]
	\centering
	\includegraphics[width=0.8\textwidth]{4.1.5.2.png}
	\caption{Confusion Matrix của LSTM + Word2Vec.}
	\label{fig:bieudo16}
\end{figure}

Sự áp đảo của True Negative (TN): Tương tự như mạng CNN, các ô True Negative ở tất cả các nhãn đều có giá trị rất cao (ví dụ: admiration đạt 4763, amusement đạt 5040). Điều này phản ánh chính xác bản chất thưa thớt của bài toán đa nhãn, nơi phần lớn các bình luận sẽ không chứa sắc thái cảm xúc đang được xét đến.

Khả năng nhận diện tốt ở các nhãn có ngữ cảnh đặc trưng: Đối với các cảm xúc có khối lượng dữ liệu huấn luyện dồi dào và cấu trúc câu rõ ràng, LSTM cho thấy khả năng nắm bắt trật tự từ vựng hiệu quả. Số ca đoán trúng (True Positive) ở các nhãn như gratitude (281), amusement (256) duy trì ở mức cao và tỷ lệ bỏ sót (False Negative) tương đối thấp.

Sự "tê liệt" trước các nhãn thiểu số: Hạn chế lớn nhất của mô hình tiếp tục bộc lộ tại các nhãn cực hiếm. Nhãn grief có số ca đoán trúng (TP) bằng 0, nhãn pride chỉ đoán trúng 2 ca. Dù kiến trúc bộ nhớ dài - ngắn của LSTM được thiết kế để học ngữ cảnh phức tạp, sự thiếu hụt trầm trọng của các mẫu dữ liệu này khiến mạng nơ-ron không đủ cơ sở để thiết lập ranh giới phân loại, dẫn đến việc mô hình gần như bỏ qua hoàn toàn các trạng thái tâm lý này.

Hiện tượng nhầm lẫn lớn ở nhãn Neutral: Đáng chú ý, tại nhãn Trung tính (Neutral), lượng dự đoán nhầm (False Positive) lên tới con số 1043 mẫu. Điều này chứng tỏ khi đối mặt với các câu bình luận lấp lửng, mơ hồ hoặc chứa các từ lóng chưa từng xuất hiện trong tập huấn luyện, LSTM vẫn gặp lúng túng và có xu hướng "đẩy" dự đoán về trạng thái an toàn là trung tính.

\begin{figure}[H]
	\centering
	\includegraphics[width=0.8\textwidth]{4.1.5.3.png}
	\caption{Biểu đồ lịch sử huấn luyện của LSTM + Word2Vec.}
	\label{fig:bieudo4.1.5.3}
\end{figure}

\textbf{Đánh giá biểu đồ huấn luyện (Learning Curve)}
\begin{itemize}
	\item Thông qua biểu đồ Training vs Validation Loss, có thể thấy mạng LSTM hội tụ rất tốt ở những Epoch đầu tiên, với Validation Loss giảm dần và chạm đáy ở khoảng Epoch 4.
	
	\item Đồng thời, đồ thị Validation Binary Accuracy cũng đạt đỉnh và duy trì ổn định ở mức trên 0.96 tại Epoch 4.
	
	\item Tuy nhiên, từ Epoch thứ 5 trở đi, hiện tượng quá khớp (Overfitting) bắt đầu bộc lộ. Training Loss tiếp tục giảm tiệm cận 0, nhưng Validation Loss lại đảo chiều tăng vọt lên mức trên 0.035. Nhờ vào cơ chế Dừng sớm (Early Stopping) được cài đặt, hệ thống đã tự động dừng huấn luyện và khôi phục bộ trọng số tối ưu nhất tại Epoch 4 để tiến hành kiểm thử
\end{itemize}

\section{So sánh}

Dựa trên kết quả thực nghiệm của 5 mô hình trên cùng một tập dữ liệu thử nghiệm, nhóm rút ra các so sánh tổng thể như sau:
\begin{itemize}
	\item \textbf{Hiệu năng tổng quát (Macro F1-Score)}: Các mô hình tuyến tính truyền thống bao gồm Hồi quy Logistic (0.43) và LinearSVC (0.43) kết hợp với biểu diễn TF-IDF bất ngờ thể hiện hiệu năng cao nhất và đồng đều nhất trên toàn hệ thống. Mô hình học tổ hợp Random Forest đạt 0.41, trong khi các mạng học sâu 1D-CNN và LSTM đạt điểm số lần lượt là 0.40 và 0.39.
	
	\item \textbf{Khả năng xử lý mất cân bằng dữ liệu}: Các mô hình áp dụng chiến lược One-Vs-Rest (như LinearSVC và Hồi quy Logistic) kiểm soát ranh giới phân loại rất tốt trên các không gian đặc trưng thưa thớt (sparse matrix). Thuật toán tối đa hóa lề (max-margin) của SVM giúp nó có khả năng "vớt vát" được độ bao phủ (Recall) cho các nhãn cực hiếm. Ngược lại, Random Forest hay mạng học sâu lại rất dễ bị tổn thương ở phần đuôi dài (long-tail) của phân bố dữ liệu, khiến kết quả dự đoán nhãn thiểu số gần như sụp đổ.
\end{itemize}

\section{Ưu nhược điểm của các mô hình}

\textbf{Hồi quy Logistic \& Support Vector Machine (LinearSVC)}:
\begin{itemize}
	\item \textbf{Ưu điểm}: Khả năng xử lý xuất sắc ma trận thưa thớt (5000 chiều) sinh ra từ phương pháp TF-IDF. Tốc độ huấn luyện nhanh, có cơ chế điều chuẩn linh hoạt, và đặc biệt tối ưu với chiến lược dò tìm siêu tham số tự động.
	
	\item \textbf{Nhược điểm}: Do đặc thù của TF-IDF là bỏ qua thứ tự từ (Bag-of-Words), các mô hình này không thể hiểu được ngữ cảnh vĩ mô của câu. Hậu quả là chúng dễ bị bối rối trước các câu lấp lửng và có xu hướng "đổ đống" dự đoán sai (False Positive) vào nhãn Neutral.
\end{itemize}

\textbf{Random Forest}:
\begin{itemize}
	\item \textbf{Ưu điểm}: Tận dụng trí tuệ tập thể (Ensemble Learning) giúp mô hình đạt độ chuẩn xác (Precision) cực kỳ cao và ổn định trên các nhãn đa số (Amusement, Joy, Anger).
	
	\item \textbf{Nhược điểm}: Hiệu năng sụt giảm nghiêm trọng đối với dữ liệu mất cân bằng. Các nhãn thiểu số (Grief, Pride) không cung cấp đủ lượng mẫu để cây quyết định phân nhánh học luật, khiến Recall chạm đáy.
\end{itemize}

\textbf{Mạng học sâu (1D-CNN \& LSTM)}:
\begin{itemize}
	\item \textbf{Ưu điểm}: Kết hợp với Word2Vec giúp mô hình thấu hiểu mối liên hệ ngữ nghĩa (semantic similarity) giữa các từ vựng. CNN bóc tách xuất sắc các cụm n-gram biểu cảm, còn LSTM xử lý tốt các câu văn có độ phụ thuộc ngữ cảnh dài.
	
	\item \textbf{Nhược điểm}: Dễ gặp hiện tượng học vẹt (Overfitting) trên tập dữ liệu mất cân bằng. Đòi hỏi tài nguyên phần cứng lớn (GPU) để tối ưu hóa hàng nghìn chu kỳ lặp và phụ thuộc rất nhiều vào các cơ chế như Early Stopping hay Threshold Tuning để bảo toàn hiệu năng thực tế.
\end{itemize}

\section{Kết quả thực nghiệm trực quan trên văn bản thực tế}

\subsection{Kịch bản thực nghiệm}

Để đánh giá và đối chứng trực quan năng lực suy luận của 05 mô hình phân loại (bao gồm cả các tiếp cận Học máy truyền thống và Học sâu), nhóm nghiên cứu xây dựng kịch bản thực nghiệm định tính (Live Test) bằng cách chọn lọc 04 câu bình luận tiếng Anh ngẫu nhiên chưa từng xuất hiện trong tập huấn luyện. Các câu văn này được thiết kế theo 04 ngữ cảnh thực tế khác nhau nhằm thách thức toàn diện các bộ lọc: (1) Câu mang tính chất cổ vũ chứa đa sắc thái phức tạp; (2) Câu chứa các từ ngữ tương phản dễ gây hiểu lầm về mặt ngữ nghĩa thuần túy (\textit{laughed so hard... ribs hurt}); (3) Câu tường minh mang đặc trưng từ vựng mạnh (\textit{Thank you so much}); và (4) Câu thể hiện trạng thái tâm lý hỗn hợp chuyển biến qua liên từ (\textit{nervous... but also excited}). Cả 04 kịch bản này sau đó sẽ được đưa qua cùng một đường ống tiền xử lý văn bản để tiến hành chạy thử nghiệm đồng thời trên cả 05 mô hình

\subsection{Kết quả thực nghiệm}

\textbf{Hồi quy Logistic}

\begin{figure}[H]
	\centering
	\includegraphics[width=0.8\textwidth]{4.5.2.1.png}
	\caption{Kết quả thực nghiệm Logistics.}
	\label{fig:bieudo4.5.2.1}
\end{figure}

Mô hình Hồi quy Logistic (OVR) thể hiện khả năng phân loại đa nhãn ổn định, nhận diện chính xác các từ khóa tường minh để kích hoạt đúng nhãn chủ đạo như \textit{Gratitude} (83.58\% ở kịch bản 3) hay cặp trạng thái hỗn hợp \textit{Excitement} (75.66\%) - \textit{Nervousness} (74.13\% ở kịch bản 4). Tuy nhiên, do hạn chế đặc trưng của học máy truyền thống là chưa tối ưu được ngữ cảnh chuỗi, mô hình bị giảm độ tự tin ở câu có cấu trúc ẩn dụ phức tạp. Minh chứng ở kịch bản 2, hai sắc thái đối lập \textit{Amusement} (56.72\%) và \textit{Sadness} (60.32\%) chỉ đạt mức điểm mấp mém vượt ngưỡng tối ưu. Dù vậy, đây vẫn là thuật toán baseline có hiệu năng thực tế rất tốt.

\textbf{Random Forest}

\begin{figure}[H]
	\centering
	\includegraphics[width=0.8\textwidth]{4.5.2.2.png}
	\caption{Kết quả thực nghiệm Random Forest.}
	\label{fig:bieudo4.5.2.2}
\end{figure}

Mô hình Random Forest thể hiện năng lực dự đoán đa nhãn chính xác và độ tự tin đồng đều trên cả 04 kịch bản nhờ cấu trúc tập hợp các cây quyết định. Thuật toán phân tách rất tốt các từ khóa mang đặc trưng mạnh để đưa ra xác suất vượt xa ngưỡng chuẩn, tiêu biểu như \textit{Optimism} (85.09\% ở kịch bản 1) hay \textit{Gratitude} (84.67\% ở kịch bản 3). Đặc biệt, ở câu văn phức tạp mang tính ẩn dụ ở kịch bản 2, Random Forest vẫn giữ được sự cân bằng và độ tự tin cao khi kích hoạt đồng thời hai nhãn đối lập là \textit{Amusement} (74.53\%) và \textit{Sadness} (66.89\%). Kết quả này cho thấy thuật toán có khả năng bao quát dữ liệu tốt, giảm thiểu hiện tượng quá khớp (overfitting) và là một đại diện học máy truyền thống cực kỳ mạnh mẽ.

\textbf{LinearSVC}

\begin{figure}[H]
	\centering
	\includegraphics[width=0.8\textwidth]{4.5.2.3.png}
	\caption{Kết quả thực nghiệm LinearSVC.}
	\label{fig:bieudo4.5.2.3}
\end{figure}

Mô hình LinearSVC (OVR) đạt hiệu năng nhận diện đa nhãn rất ấn tượng, thể hiện qua khoảng cách hình học Z đến đường ranh giới quyết định có độ phân tách lớn ở các nhãn chủ đạo như \textit{Gratitude} (Z = 1.3457 ở kịch bản 4) hay \textit{Excitement} (Z = 1.1642 ở kịch bản 5). Điểm cộng lớn của thuật toán này là khả năng tối ưu hóa ranh giới tuyến tính rất dứt khoát, giúp cô lập tốt các đặc trưng từ vựng mạnh. Tuy nhiên, tương tự như Hồi quy Logistic, do bản chất không học chuỗi ngữ cảnh, độ tự tin của LinearSVC bị kéo thấp đáng kể khi gặp các câu ẩn dụ ở kịch bản 2 và 3, khiến các nhãn đối lập như \textit{Amusement} (Z = 0.4280) hay \textit{Disappointment} (Z = 0.2704) rơi vào vùng mấp mém sát đường ranh giới kích hoạt.

\textbf{CNN}

\begin{figure}[H]
	\centering
	\includegraphics[width=0.8\textwidth]{4.5.2.4.png}
	\caption{Kết quả thực nghiệm CNN.}
	\label{fig:bieudo4.5.2.4}
\end{figure}

Mô hình mạng tích chập 1D CNN thể hiện khả năng trích xuất đặc trưng cục bộ rất nhạy bén nhờ cơ chế trượt của các bộ lọc (kernels). Sức mạnh này minh chứng rõ nhất ở kịch bản 3 khi cô lập và bắt trúng nhãn \textit{Gratitude} với độ tự tin gần như tuyệt đối (98.84\%), cao nhất trong số các mô hình thử nghiệm. Tuy nhiên, do cấu trúc CNN tập trung tóm lấy các tổ hợp n-gram mang tín hiệu cảm xúc mạnh nhất (qua lớp Max Pooling) thay vì duy trì toàn bộ ngữ cảnh chuỗi tuyến tính dài, mô hình có xu hướng bỏ sót các sắc thái mờ nhạt hơn trong câu phức. Cụ thể ở kịch bản 1, nhãn \textit{Gratitude} bị rơi rụng, hay ở kịch bản 4 mô hình chỉ giữ lại nhãn \textit{Excitement} (20.64\%) mà bỏ qua sự lo lắng (\textit{Nervousness}). Dù vậy, với các câu ngắn mang từ khóa đặc trưng như kịch bản 2 (\textit{Amusement} - 64.23\%), CNN vẫn chứng minh tốc độ phản hồi và năng lực nhận diện cục bộ cực kỳ đáng tin cậy.

\textbf{LSTM}

\begin{figure}[H]
	\centering
	\includegraphics[width=0.8\textwidth]{4.5.2.5.png}
	\caption{Kết quả thực nghiệm LSTM.}
	\label{fig:bieudo4.5.2.5}
\end{figure}

Mô hình mạng bộ nhớ dài-ngắn hạn LSTM (Bi-LSTM) thể hiện khả năng bao quát ngữ cảnh chuỗi rất tốt nhờ cơ chế các cổng nhớ (gates) duy trì thông tin tuyến tính. Đặc tính này giúp mô hình nhận diện mượt mà các tổ hợp câu dài mang sắc thái hỗn hợp, kích hoạt chính xác bộ ba nhãn \textit{Admiration} (58.43\%), \textit{Caring} (30.20\%), \textit{Optimism} (54.60\%) ở kịch bản 1, hay cô lập xuất sắc nhãn \textit{Gratitude} với độ tự tin cao (87.17\%) ở kịch bản 3. Tuy nhiên, so với các thuật toán học máy truyền thống hoặc cấu trúc CNN tập trung tóm sững đặc trưng mạnh cục bộ, mạng LSTM đôi khi có xu hướng phân phối xác suất 'phẳng' hơn khi đối mặt với các câu ẩn dụ hoặc chuyển biến trạng thái đột ngột. Hệ quả là ở kịch bản 2, mô hình chỉ giữ lại nhãn \textit{Amusement} (54.49\%) mà bỏ sót sắc thái tiêu cực đi kèm, và ở kịch bản 4 lại kích hoạt nhãn \textit{Joy} (45.50\%) thay vì sự lo lắng (\textit{Nervousness}). Dù vậy, LSTM vẫn chứng minh được ưu thế lớn về khả năng duy trì mạch ngữ nghĩa ổn định xuyên suốt văn bản.

\subsection{So sánh kết quả thực nghiệm 5 mô hình}

Dựa trên kết quả chạy thử nghiệm trực quan (Live Test) của cả 05 mô hình phân loại đa nhãn cảm xúc (Hồi quy Logistic, Random Forest, LinearSVC, CNN, và LSTM), nhóm nghiên cứu đưa ra các đánh giá đối chứng cụ thể như sau:
\begin{itemize}
	\item \textbf{Về năng lực nhận diện cấu trúc câu tường minh}: Tất cả 05 mô hình đều hoàn thành xuất sắc nhiệm vụ khi cô lập chính xác duy nhất nhãn \textit{Gratitude}. Trong đó, mô hình mạng tích chập CNN đạt độ tự tin cao nhất với 98.84\%, bám sát phía sau là LSTM (87.17\%), Random Forest (84.67\%), và Hồi quy Logistic (83.58\%). Điều này chứng minh khi văn bản chứa các từ khóa đặc trưng mạnh ("\textit{thank you}"), ranh giới phân loại của cả phương pháp truyền thống lẫn học sâu đều hoạt động rất tường minh.
	
	\item \textbf{Về khả năng xử lý câu chứa từ ngữ tương phản/ẩn dụ}: Đây là kịch bản chứng kiến sự phân hóa rõ rệt nhất về mặt kiến trúc thuật toán. Mô hình Random Forest thể hiện sự cân bằng đáng kinh ngạc khi kích hoạt song song hai nhãn đối lập \textit{Amusement} (74.53\%) và \textit{Sadness} (66.89\%) với độ tự tin cao. Mô hình Hồi quy Logistic và LinearSVC cũng nhận diện được hai sắc thái này nhưng điểm số rơi vào vùng mập mém sát ngưỡng kích hoạt. Ngược lại, các mạng học sâu (CNN và LSTM) do cơ chế tiền xử lý chặt chẽ và xu hướng tập trung vào các đặc trưng tối ưu hóa phân phối, nên chỉ kích hoạt nhãn chủ đạo \textit{Amusement} mà chấp nhận bỏ sót sắc thái tiêu cực đi kèm.
	
	\item \textbf{Về khả năng bao quát ngữ cảnh đa nhãn phức tạp}: Mô hình mạng LSTM chứng minh được ưu thế của cơ chế các cổng nhớ khi duy trì tốt mạch thông tin tuyến tính, giúp kích hoạt trọn vẹn bộ ba nhãn (Admiration, Caring, Optimism) ở câu chúc và bóc tách được các trạng thái chuyển biến tâm lý phức tạp. Trong khi đó, CNN do đặc tính tích chập và rút gọn thông tin cục bộ (Max Pooling) nên có xu hướng chỉ giữ lại các cụm n-gram có tín hiệu mạnh nhất, dẫn đến việc giải quyết bài toán đa nhãn ở kịch bản phức tạp bị thu hẹp hoặc sót nhãn so với các đối thủ còn lại.
\end{itemize}

Không có mô hình nào áp đảo tuyệt đối trên mọi kịch bản định tính. Các mô hình Học máy truyền thống (Random Forest, LinearSVC, Logistic Regression) thể hiện tư duy phân phối xác suất rất đều và nhạy với phân loại đa nhãn trên từ khóa. Trong khi đó, các mô hình Học sâu (CNN, LSTM) lại sở hữu những "đỉnh nhọn" về độ tự tin cực cao (gần như tuyệt đối) ở các phân đoạn ngữ cảnh cốt lõi được huấn luyện kỹ. 
% Chương 5

\chapter{KẾT LUẬN VÀ HƯỚNG PHÁT TRIỂN} % Tên của chương

\label{Chapter5} % Để trích dẫn chương này ở chỗ nào đó trong bài, hãy sử dụng lệnh \ref{Chapter5}

%----------------------------------------------------------------------------------------

\section{Kết quả đạt được}

Đề tài đã hoàn thành xuất sắc các mục tiêu cốt lõi được đặt ra ban đầu, bao gồm việc nghiên cứu, xây dựng và đánh giá toàn diện hệ thống phân loại đa nhãn cảm xúc từ dữ liệu văn bản mạng xã hội. Bằng việc khai thác bộ dữ liệu chuẩn GoEmotions với 28 trạng thái tâm lý phức tạp, nhóm nghiên cứu đã triển khai thành công một luồng xử lý ngôn ngữ tự nhiên (NLP) khép kín. Quá trình này bắt đầu từ việc làm sạch dữ liệu thô, chuẩn hóa ngôn từ mạng xã hội, xử lý cấu trúc phủ định, cho đến bước trích xuất đặc trưng văn bản thông qua hai phương pháp là thống kê TF-IDF và kỹ thuật nhúng từ Word2Vec. Qua quá trình thực nghiệm và so sánh chéo 5 thuật toán học máy đại diện bao gồm Hồi quy Logistic, Random Forest, LinearSVC, 1D-CNN và LSTM, kết quả cho thấy các thuật toán tuyến tính truyền thống kết hợp với ma trận thưa TF-IDF (đặc biệt là LinearSVC và Logistic Regression) mang lại hiệu quả khái quát hóa ấn tượng nhất với điểm Macro F1-Score đạt 0.43. Khả năng phân định ranh giới rõ ràng thông qua chiến lược One-Vs-Rest đã giúp các mô hình này xử lý tốt đặc tính văn bản ngắn và vượt qua rào cản về sự mất cân bằng cực đoan dạng đuôi dài (long-tail) của dữ liệu. Trong khi đó, các mô hình học sâu (CNN, LSTM) dù sở hữu khả năng nắm bắt ngữ nghĩa tốt nhưng lại bộc lộ hạn chế ở các nhãn thiểu số do thiếu hụt lượng mẫu huấn luyện cần thiết, đồng thời đòi hỏi sự tinh chỉnh phức tạp về cơ chế chống quá khớp và dò tìm ngưỡng tối ưu tự động (Threshold Tuning).

\section{Hướng phát triển trong tương lai}

Mặc dù hệ thống đã đạt được những kết quả khả quan trong việc bóc tách và nhận diện các cụm cảm xúc phổ biến, song năng lực dự đoán đối với các trạng thái tâm lý ngầm ẩn, hiếm gặp (như \textit{grief}, \textit{pride}, \textit{relief}) vẫn là một thách thức lớn cần được giải quyết. Trong tương lai, định hướng quan trọng nhất của đề tài là áp dụng các kỹ thuật cân bằng dữ liệu nâng cao, điển hình như phương pháp sinh mẫu thiểu số (SMOTE) hoặc tăng cường dữ liệu (Data Augmentation), nhằm cung cấp đủ không gian ngữ cảnh cho mô hình học tập ở những nhãn hiếm. Bên cạnh đó, việc nâng cấp hệ thống trích xuất đặc trưng từ Word2Vec lên các kiến trúc Mô hình ngôn ngữ lớn (Large Language Models) tiên tiến dựa trên nền tảng Transformer như BERT hoặc RoBERTa sẽ là một bước đột phá. Các mô hình này có khả năng "thấu cảm" sâu sắc hơn cấu trúc ngữ pháp, nắm bắt chính xác hiện tượng đa nghĩa và những sắc thái tinh vi của con người trên không gian mạng. Cuối cùng, nhóm hướng tới việc tối ưu hóa và đóng gói mô hình thành một hệ thống API hoàn chỉnh, mở ra cơ hội ứng dụng thực tiễn trong các lĩnh vực như phân tích rủi ro tâm lý người dùng diễn đàn, kiểm duyệt nội dung tự động hoặc hỗ trợ phân tích phản hồi trong dịch vụ chăm sóc khách hàng

\section{Tài liệu mã nguồn mở}

Nhằm đảm bảo tính minh bạch, tính lặp lại của các thực nghiệm khoa học, cũng như tạo tài liệu tham khảo cho các nghiên cứu phát triển tiếp theo, toàn bộ mã nguồn lập trình, các kịch bản tiền xử lý dữ liệu (notebooks), ma trận trọng số của các mô hình đã huấn luyện và tài liệu liên quan đến đề tài đều được nhóm lưu trữ và công khai hoàn toàn.
Chi tiết hệ thống và mã nguồn dự án được lưu trữ tại kho mã nguồn GitHub của nhóm:

Đường dẫn truy cập: https://github.com/trongtung2005/Machine\_Learning\_2526II\_Nhom6 

%----------------------------------------------------------------------------------------
%	(KHÔNG CHỈNH SỬA PHẦN NÀY)
%
%	PHẦN 12: TÀI LIỆU THAM KHẢO
%----------------------------------------------------------------------------------------

\begin{spacing}{1.15}
	\printbibliography[heading=bibintoc, title=Tài liệu tham khảo] % In ra tài liệu tham khảo
\end{spacing}

%----------------------------------------------------------------------------------------
%	PHẦN 13: PHỤ LỤC (THESIS CONTENT - APPENDICES)
%----------------------------------------------------------------------------------------

\appendix % Nói với LaTeX rằng những chương về sau được tính là phụ lục

% Hãy thêm những phụ lục (appendix) của khóa luận/tiểu luận vào thư mục Appendices
% Hãy bỏ chú thích những dòng nếu bạn đã bổ sung những phụ lục vào

% Phụ lục A

\chapter{Định nghĩa các nhãn cảm xúc} % Tên của phụ lục

\label{AppendixA} % Để trích dẫn chương này ở chỗ nào đó trong bài, hãy sử dụng lệnh \ref{AppendixA} 

%----------------------------------------------------------------------------------------

Hệ thống nhãn cảm xúc trong bộ dữ liệu GoEmotions bao gồm 27 danh mục sắc thái cảm xúc chi tiết và 1 nhãn trung tính, được định nghĩa cụ thể như sau:

\begin{enumerate}
	\item \textbf{admiration} \twemoji{1f44f}: NGƯỠNG MỘ (Thể hiện việc tìm thấy một điều gì đó ấn tượng hoặc đáng trân trọng, kính phục).
	
	\item \textbf{amusement} \twemoji{1f602}: THÍCH THÚ (Thể hiện việc tìm thấy một điều gì đó hài hước, vui nhộn hoặc mang tính giải trí).
	
	\item \textbf{anger} \twemoji{1f621}: GIẬN DỮ (Cảm giác không hài lòng mạnh mẽ, bất bình hoặc mang tính phản kháng, đối địch).
	
	\item \textbf{annoyance} \twemoji{1f612}: KHÓ CHỊU (Trạng thái tức giận nhẹ, bực mình hoặc kích động không thoải mái).
	
	\item \textbf{approval} \twemoji{1f44d}: TÁN THÀNH (Việc sở hữu hoặc biểu đạt một quan điểm, ý kiến thuận chiều, ủng hộ).
	
	\item \textbf{caring} \twemoji{1f970}: QUAN TÂM (Biểu lộ sự tử tế, lòng trắc ẩn và sự lo lắng, chú ý dành cho người khác).
	
	\item \textbf{confusion} \twemoji{1f928}: BỐI RỐI (Sự thiếu thấu hiểu, mơ hồ hoặc trạng thái không chắc chắn về một vấn đề).
	
	\item \textbf{curiosity} \twemoji{1f914}: TÒ MÒ (Mong muốn mạnh mẽ được biết, được khám phá hoặc học hỏi một điều gì đó).
	
	\item \textbf{desire} \twemoji{1f97a}: KHAO KHÁT (Cảm giác mạnh mẽ muốn có được một thứ gì đó hoặc mong ước một sự việc cụ thể xảy ra).
	
	\item \textbf{disappointment} \twemoji{1f61e}: THẤT VỌNG (Sắc thái buồn bã hoặc không hài lòng gây ra bởi việc không đạt được kỳ vọng hay hy vọng).
	
	\item \textbf{disapproval} \twemoji{1f44e}: PHẢN ĐỐI (Việc sở hữu hoặc biểu đạt một quan điểm, ý kiến không thuận chiều, bất đồng).
	
	\item \textbf{disgust} \twemoji{1f92e}: GHÊ TỞM (Sự phẫn nộ hoặc phản đối mạnh mẽ được khơi dậy bởi một điều gì đó khó chịu hoặc phản cảm).
	
	\item \textbf{embarrassment} \twemoji{1f633}: NGƯỢNG NGÙNG (Trạng thái tự ý thức tiêu cực, xấu hổ, ngượng nghịu hoặc lúng túng trong ngữ cảnh xã hội).
	
	\item \textbf{excitement} \twemoji{1f929}: HÀO HỨNG (Cảm giác nhiệt huyết lớn, mong đợi mãnh liệt và tràn đầy năng lượng tích cực).
	
	\item \textbf{fear} \twemoji{1f628}: SỢ HÃI (Trạng thái hoảng sợ, lo lắng sâu sắc hoặc e ngại trước một mối nguy hại).
	
	\item \textbf{gratitude} \twemoji{1f64f}: BIẾT ƠN (Cảm giác biết ơn, trân trọng và đánh giá cao sự giúp đỡ hoặc một điều tốt đẹp nhận được).
	
	\item \textbf{grief} \twemoji{1f62d}: ĐAU BUỒN (Nỗi đau khổ dữ dội, tột cùng, đặc biệt là nỗi đau gây ra bởi sự ra đi của một ai đó).
	
	\item \textbf{joy} \twemoji{1f600}: VUI SƯỚNG (Cảm giác cực kỳ hạnh phúc, thỏa mãn và tràn ngập niềm vui).
	
	\item \textbf{love} \twemoji{2764}: YÊU THƯƠNG (Cảm xúc tích cực mạnh mẽ thể hiện sự gắn kết, quý mến và tình cảm sâu sắc).
	
	\item \textbf{nervousness} \twemoji{1f62c}: LO LẮNG (Trạng thái bồn chồn, e sợ, căng thẳng hoặc bất an trước một sự kiện).
	
	\item \textbf{optimism} \twemoji{1f91e}: LẠC QUAN (Sự hy vọng và niềm tin tưởng mạnh mẽ vào tương lai hoặc sự thành công của một công việc).
	
	\item \textbf{pride} \twemoji{1f60e}: TỰ HÀO (Niềm vui lớn hoặc sự thỏa mãn sâu sắc đến từ thành tựu của chính mình hoặc của những người gắn bó).
	
	\item \textbf{realization} \twemoji{1f4a1}: NHẬN RA (Trạng thái bắt đầu nhận thức, thấu hiểu hoặc đột nhiên thức tỉnh về một sự thật nào đó).
	
	\item \textbf{relief} \twemoji{1f60c}: NHẸ NHÕM (Sự trấn tĩnh, an tâm và thư giãn sau khi được giải phóng khỏi trạng thái lo âu hoặc đau khổ).
	
	\item \textbf{remorse} \twemoji{1f625}: HỐI HẬN (Cảm giác hối tiếc sâu sắc hoặc dằn vặt, có lỗi về một hành vi sai lầm đã qua).
	
	\item \textbf{sadness} \twemoji{1f622}: BUỒN BÃ (Nỗi đau đớn về mặt cảm xúc, tâm trạng sầu não, chán nản).
	
	\item \textbf{surprise} \twemoji{1f62e}: NGẠC NHIÊN (Cảm giác kinh ngạc, sửng sốt hoặc giật mình gây ra bởi một điều gì đó hoàn toàn bất ngờ).
	
	\item \textbf{neutral} \twemoji{1f610}: TRUNG TÍNH (Văn bản thuần túy thông tin thực tế, không chứa bất kỳ sắc thái biểu cảm hay cảm xúc nào).
\end{enumerate}
% Phụ lục B

\chapter{Thống kê phân phối số lượng nhãn trên mỗi mẫu dữ liệu} % Tên của phụ lục

\label{AppendixB} % Để trích dẫn chương này ở chỗ nào đó trong bài, hãy sử dụng lệnh \ref{AppendixB} 

%----------------------------------------------------------------------------------------

Nhãn cảm xúc được coi là 'chỉ được chọn bởi một người gán nhãn duy nhất' khi trong một nhóm từ 3 đến 5 người cùng đọc một câu văn, chỉ có đúng một người lựa chọn nhãn đó còn những người khác thì không. Đây là những nhãn mang tính chủ quan cá nhân cao, dễ gây nhiễu nên đã được thuật toán lọc bỏ (chuyển từ pre-filter sang post-filter) để đảm bảo tính đồng thuận và giúp bộ dữ liệu sạch hơn trước khi đưa vào huấn luyện mô hình

\begin{figure}[H]
	\centering
	\includegraphics[width=0.8\textwidth]{B.png}
	\caption{Thống kê số lượng nhãn cảm xúc trên mỗi mẫu dữ liệu trước và sau khi lọc bỏ các nhãn chỉ được chọn bởi một người gắn nhãn duy nhất.}
	\label{fig:B}
\end{figure}
% Phụ lục C

\chapter{Bảng phân công công việc} % Tên của phụ lục

\label{AppendixC} % Để trích dẫn chương này ở chỗ nào đó trong bài, hãy sử dụng lệnh \ref{AppendixC}

%----------------------------------------------------------------------------------------

\begin{table}[htbp]
	\centering
	\caption{Bảng phân công công việc và đóng góp của các thành viên Nhóm 6.}
	\label{tab:phancông_nhom6}
	\footnotesize % Giảm cỡ chữ xuống một chút để khít hoàn toàn vào trang giấy, không bị tràn viền
	\begin{tabularx}{\textwidth}{|c|p{2.8cm}|c|X|c|c|c|}
		\hline
		\textbf{STT} & \textbf{Họ và tên} & \textbf{Mã SV} & \textbf{Phân công nhiệm vụ} & \begin{tabular}[c]{@{}c@{}}\textbf{Mức độ}\\\textbf{hoàn thành}\end{tabular} & \begin{tabular}[c]{@{}c@{}}\textbf{Trọng số}\\\textbf{cá nhân}\end{tabular} & \textbf{Chữ ký} \\ \hline
		1 & Nguyễn Doãn Bình & 23001582 & Nội dung quyển báo cáo \newline Nội dung Slide & 80\% & 24.50\% & Bình \\ \hline
		2 & Phan Thành Đạt & 23001595 & Nội dung quyển báo cáo \newline Nội dung Slide \newline Thuyết trình & 80\% & 24.50\% & Đạt \\ \hline
		3 & Tô Quang Hân & 23001605 & Nội dung quyển báo cáo \newline Nội dung Slide \newline Code mô hình RandomForest và LinearSVC & 87\% & 25.50\% & Hân \\ \hline
		4 & Nguyễn Trọng Tùng & 23001638 & Nội dung quyển báo cáo \newline Giao diện và nội dung slide \newline Code mô hình Logistics, CNN, LSTM \newline Tinh chỉnh lại tất cả các mô hình & 90\% & 25.50\% & Tùng \\ \hline
	\end{tabularx}
\end{table}

%----------------------------------------------------------------------------------------

\end{document} 
%----------------------------------------------------------------------------------------------------------------------------------------
